% main_v15.tex = main_v14 with sec_7_verification_v12 (Table tab:finite2 entry of the adversarial search; "checked in Lean"),
% the Data availability sentence on the recorded witnesses of Theorem G, and the software concept DOI (all versions) instead of the version 1.0.0 DOI.
% main_v14.tex = main_v13 with the repository URL and the Zenodo DOI of the software record in Data and code availability
% (identifiers reserved 2026-10-13). No other change.
% main_v13.tex = main_v12 with the Data and code availability paragraph matched to the released checks/ and evidence/,
% sec_7_verification_v11, app_A_mixing_v2 (layout only) and bib_v2 ([FukuiFFA] paper DOI).
% main_v12.tex -- public (named) build, C9. = main_v11 with the author block and the AI statement aligned with the
% author's approved series data (weber_general_n MoC manuscript: ORCID inside \address, no postal address, no \urladdr).
% Shared with main_v7_tosc_anon.tex: macros_v1, abstract_v1, all section files, app_A_mixing_v1, bib_v1.
% (v8..v10 were intermediate: table split, float placement; see c9/notes/resp/P9-08.md.)
% Author data: c9/paper/AUTHOR_CHECK_C9.md (items marked there need the author's confirmation).
\documentclass[11pt]{amsart}
\usepackage{amsmath,amssymb,amsthm}
\usepackage{booktabs}
\usepackage{url}
\usepackage[hidelinks]{hyperref}
\hypersetup{
  pdftitle={Carry-aware differential bounds and certified correction of local zero sums},
  pdfauthor={Hiroki Fukui},
  pdfsubject={Galois rings, APN permutations, key-averaged differential bounds},
  pdfkeywords={APN permutation, Galois ring, differential uniformity, Markov cipher, key-averaged differential probability,
    contraction, zero sum, multiplicative complexity}
}

\newtheorem{theorem}{Theorem}[section]
\newtheorem{lemma}[theorem]{Lemma}
\newtheorem{proposition}[theorem]{Proposition}
\newtheorem{corollary}[theorem]{Corollary}
\theoremstyle{definition}
\newtheorem{definition}[theorem]{Definition}
\newtheorem{assumption}[theorem]{Assumption}
\newtheorem{remark}[theorem]{Remark}
\newtheorem{example}[theorem]{Example}

% macros_v1.tex -- C9 paper lane: notation macros shared by main_v7.tex and main_v7_tosc_anon.tex
% (identical to the \newcommand block of main_v6.tex).
\newcommand{\Fq}{\mathbb{F}_q}
\newcommand{\FF}{\mathbb{F}_2}
\newcommand{\GR}{\mathrm{GR}}
\newcommand{\Tr}{\mathrm{Tr}}
\newcommand{\wt}{\mathrm{wt}}
\newcommand{\Img}{\mathrm{Im}\,}
\newcommand{\DP}{\mathrm{DP}}
\newcommand{\EDP}{\mathrm{EDP}}
\newcommand{\E}{\mathbb{E}}
\newcommand{\cZ}{\mathcal{Z}}
\newcommand{\cE}{\mathcal{E}}
\newcommand{\cG}{\mathcal{G}}
\newcommand{\cK}{\mathcal{K}}
\newcommand{\RF}{\mathrm{RF}}
\newcommand{\RP}{\mathrm{RP}}


\title[Carry-aware differential bounds and local zero sums]{Carry-aware differential bounds and certified correction of local zero sums}
\author{Hiroki Fukui}
\address{Research Institute of Criminal Psychiatry / Sex Offender Medical Center, Tokyo, Japan; and Department of Neuropsychiatry, Kyoto University, Kyoto, Japan. ORCID: 0009-0008-7122-522X}
\email{fukui@somec.org}
\subjclass[2020]{Primary 94A60; Secondary 11T71, 11T06, 13M10, 60J10}
\keywords{APN permutation, Galois ring, differential uniformity, Markov cipher, key-averaged differential probability,
contraction, zero sum, multiplicative complexity}

\setlength{\emergencystretch}{3em}
\begin{document}

\begin{abstract}
% abstract_v1.tex -- C9 paper lane (P9-01): abstract shared by the public and the anonymous build.
% Exponent 1/2 restricted to p(x)=x^{q-2}, odd m, d(a)=1+(a/rho)^4+(a/rho)^2, rho!=0 (finding C8-F02); the general-d
% repair statement keeps "every admissible d"; q=2^m and N=q^2 are stated.
We study a two-digit coordinate model of maps on the Galois ring $\GR(4,m)$, with $q=2^m$ and $N=q^2=|\GR(4,m)|$, in which the
ring addition carries from the low digit into the high digit. For maps
$F(a,b)=(p(a),\,h(a)\oplus d(a)b)$ with $p$ an APN permutation, $d$ a nowhere-vanishing
$\FF$-affine map whose linear part has a one-dimensional kernel, and $h$ arbitrary, we prove
direction-wise differential counts for $B=J\circ F$, a key-averaged two-block bound of order $1/N$, and a
collision-energy bound, all uniform in $h$. In contrast, a local integral property of the same
two-block map depends on $h$; for the inverse map and every admissible $d$ we give a linear
certificate on $h$ that excludes a family of local zero sums, and an explicit correction of Boolean
degree two of every $h$ that fails it; degree two is minimal for this local problem. For the averaged
alternating two-word arrangement we prove a one-word contraction of order $q^{-(1-\frac12\log_2 3)}$, uniform in $h$ and in the
admissible data, and deduce that a fixed number of steps brings the averaged two-word difference distribution to the scale $1/N$.
For the inverse map with $d(a)=1+(a/\rho)^4+(a/\rho)^2$, where $m$ is odd and $\rho\ne0$, the contraction improves to the order
$q^{-1/2}$, which is optimal since no admissible data give a faster order; for other $d$ our constant grows with the linearized
degree of $d$.
We then record what these statements do and do not imply for an independent-key composition:
the averaged multi-word bound coexists with a candidate-existence distinguisher with data linear
in the word size. No security claim for a concrete cipher is made.

\end{abstract}

\maketitle
\tableofcontents

% sec_1_intro_v6.tex -- C9 paper lane: layout fix of sec_1_intro_v5.tex (Table tab:main was 10.4pt too wide; column
% separation 3pt, widths adjusted). Text identical to v5.
% sec_1_intro_v5.tex -- C9 paper lane (P9-01, P9-09, P9-10, P9-14). Changes from sec_1_intro_v4.tex (C8):
% (1) contraction paragraph: exponent 1/2 restricted to p(x)=x^{q-2}, odd m, d=1+K0(x/rho), K0=X^4+X^2, rho!=0;
%     general linearized degree e shown with its constant 2^{e+1}+5 (finding C8-F02);
% (2) new summary table of the main statements (family / all h / rho / exponent or constant / steps);
% (3) relation to the two earlier papers written in the third person, with the result numbers of the
%     submitted versions (derivation recorded in c9/lit/P9-14_lit.md); internal comments removed.
\section{Introduction}\label{sec:intro}

\subsection{Two kinds of statements about the same map}
Let $R=\GR(4,m)$ be the Galois ring of characteristic $4$ and residue field $\Fq$, $q=2^m$, and put $N=q^2=|R|$.
Writing $x\in R$ in base $2$ along a fixed basis gives two digits $x=(a,b)\in\Fq^2$; ring addition
is then the digit-wise exclusive or, plus a carry $a\star u$ from the low into the high digit
(Lemma~\ref{lem:add}). A polynomial map of $R$ has, in these coordinates, the shape
\[
  F(a,b)=\bigl(p(a),\,h(a)\oplus d(a)\,b\bigr)
\]
(Proposition~\ref{prop:realize}); the swap $J(a,b)=(b,a)$ turns $F$ into a map $B=J\circ F$ whose
low digit depends on $b$. The present paper separates two kinds of statements about $B$.

\emph{Differential statements are uniform in $h$.} Under an assumption on $(p,d)$ only
(Assumption~\ref{ass:PD}), the differential uniformity of $B$ with respect to ring subtraction is
exactly $q$ for every $h$ (Theorem~\ref{thm:eps}); the key-averaged differential probability of a
two-block local map is at most $4/N$ when the kernel direction $\rho$ of $d$ has weight one, and
at most $(6+6/q)/N$ in general, again for every $h$ (Theorem~\ref{thm:W2}); and a collision energy
is bounded by an absolute constant (Theorem~\ref{thm:UE}). The proofs are direction-wise counts
that keep the carry visible; no change of basis is used.

\emph{Contraction is uniform in $h$.} For the key-averaged alternating two-word arrangement, the one-word fluctuation
operator of $B$ has norm $\tau$ with $\tau^2\le5(3/4)^m-3/q+1/q^2$ for all PD data and all $h$ (Theorem~\ref{thm:R1}), so a
fixed number of steps brings the difference distribution to the scale $1/N$ (Theorem~\ref{thm:steps}). For one explicit family,
namely $p(x)=x^{q-2}$, $m$ odd, and $d=1\oplus K_0(x/\rho)$ with $K_0(X)=X^4+X^2$ and any $\rho\ne0$, one has $\tau^2<13.5/q$ for every $h$
(Corollary~\ref{cor:half}); since one block of the operator has norm exactly $((q/2-1)/q^2)^{1/2}$ for all PD data
(Proposition~\ref{prop:Zexact}), the exponent $1/2$ is optimal for this family. For the inverse map with another admissible $d=\varepsilon\oplus K$,
$K=\sum_{j=0}^ek_jX^{2^j}$, $k_e\ne0$, the same argument gives only $\tau^2\le(2^{e+1}+5)/q+(q/2-1)/q^2$
(Proposition~\ref{prop:Minv}, Corollary~\ref{cor:half}), whose constant grows with $e$; the exponent $1/2$ uniformly over all PD
data, or over all admissible $d$ for the inverse map, is not proved (Remark~\ref{rem:exp}).

\emph{Integral statements depend on $h$.} The ring sum of the low digit of the same two-block map
over a coset $a+2R$ is a sum of a non-affine image (Section~\ref{sec:repair}); whether it vanishes
identically depends on $h$. For $p(x)=x^{q-2}$ and every admissible $d$ we give an $\Fq$-linear functional
$\ell_d$ on the space of all $h$ whose non-vanishing excludes the local zero sum for every local key, and an
explicit correction $g_d$ of Boolean degree two, whose multiplicative complexity on its own is exactly $m-1$, for every $h$ with
$\ell_d(h)=0$; affine $h$ never suffice. The differential statements are
untouched by this correction because they never depended on $h$.

The third part records the consequences, and the limits, for an independent-key composition with
three mixing layers and local lengths $(2,2,1,1)$: a multi-word averaged bound
(Theorem~\ref{thm:P4}) holds simultaneously with a candidate-existence distinguisher whose data
is linear in $N$ (Theorem~\ref{thm:V}). We use this to state precisely why a small averaged
differential bound is not a PRP-security statement.

Table~\ref{tab:main} lists the quantifiers of the main statements. ``PD'' means all $(p,d)$ under Assumption~\ref{ass:PD}.
\begin{table}[ht]\footnotesize\setlength{\tabcolsep}{3pt}
\caption{Main statements and their quantifiers.}\label{tab:main}
\begin{tabular}{p{0.15\textwidth}p{0.25\textwidth}p{0.06\textwidth}p{0.11\textwidth}p{0.33\textwidth}}
\toprule
statement & family & all $h$ & $\rho$ & exponent / constant / steps\\
\midrule
Thm.~\ref{thm:eps} & PD & yes & any & $\Delta(B)=q$, one step\\
Thm.~\ref{thm:W2} & PD, odd $m\ge5$ & yes & $\wt(\rho)=1$; any & $4/N$; $(6+6/q)/N$ (two blocks)\\
Thms.~\ref{thm:cert}, \ref{thm:repair} & $p=x^{q-2}$, every admissible $d$, $m\ge3$ & yes$^\ast$ & any & certificate $\ell_d$; correction $g_d$\\
Thm.~\ref{thm:R1} & PD, $m\ge3$ & yes & any & $\tau\le\sqrt5\,q^{-\gamma}$, $\gamma=1-\frac12\log_23$\\
Prop.~\ref{prop:Zexact} & PD, $m\ge3$ & yes & any & $\tau\ge((q/2-1)/q^2)^{1/2}$: no exponent $>1/2$\\
Cor.~\ref{cor:half} & $p=x^{q-2}$, $d=1\oplus K_0(x/\rho)$, odd $m$ & yes & any $\rho\ne0$ & $\tau^2<13.5/q$: exponent $1/2$ (optimal)\\
Cor.~\ref{cor:half} (general) & $p=x^{q-2}$, odd $m$, $\deg_2K=e$ & yes & any & $\tau^2\le(2^{e+1}+5)/q+(q/2-1)/q^2$\\
Thm.~\ref{thm:steps} & PD, $m\ge5$, A1 & yes & any & $\|D_{14}\|_2\le1\,600\,000/N$; zero word after $15$ steps\\
Cor.~\ref{cor:inv13} & $p=x^{q-2}$, $d=1\oplus K_0(x/\rho)$, odd $m\ge5$, A1 & yes & any $\rho\ne0$ & $\|D_{13}\|_2\le24\,405/N$; zero word after $14$ steps\\
Thm.~\ref{thm:P4} & PD, $q\ge32$, $w\ge2$, $(2,2,1,1)$ & yes & $\wt(\rho)=1$; any & $(4/N)^w$; $((6+6/q)/N)^w$\\
\bottomrule
\end{tabular}

\smallskip
{\footnotesize $^\ast$The certificate is a condition on $h$; the correction $h\mapsto h+g_d$ applies to every $h$ that fails it.
$\deg_2K=e$ means $K=\sum_{j\le e}k_jX^{2^j}$ with $k_e\ne0$. All bounds concern key-averaged transitions.}
\end{table}

\subsection{Relation to two earlier papers}
The papers \cite{FukuiIT} and \cite{FukuiFFA} concern standard representatives of APN permutations and their
Galois-ring lifts. The table lists the results of those papers that are used or recalled here, and how the present paper differs.
\begin{center}\small
\begin{tabular}{p{0.08\textwidth}p{0.42\textwidth}p{0.38\textwidth}}
\toprule
Paper & Result recalled & Present paper\\
\midrule
\cite{FukuiIT} & Every quadratic APN permutation in odd dimension $m\ge5$ is linearly equivalent to an optimal
one (\cite[Main Theorem]{FukuiIT}; existence). & Used only as background; no result here depends on it.\\
\cite{FukuiFFA} & Explicit reduced representatives $P_{m,r}$, permutation and APN, with derivative
$P'=1+X^{2^{r-2}}+X^{2^{m-2}}$, whose linear part has kernel $\{0,\rho\}$, $\rho=1$ (``derivative-optimal'')
and product formula $P_{m,r}+1=DE$ (\cite[Thm.~3.1, Def.~2.1]{FukuiFFA}); integer product form (\cite[Prop.~5.1]{FukuiFFA});
coefficient lifts are permutations (\cite[Prop.~6.1]{FukuiFFA}); $\delta_F\le 2$ for unit directions
$\bar\alpha\notin\{0,\rho\}$ (\cite[Thm.~6.2(1)]{FukuiFFA}). & The data $(p,d,h)$ is arbitrary subject to
Assumption~\ref{ass:PD}; $d$ need not be $p'$, $p$ need not be quadratic; the $k=2$ direction bound
is reproved (Proposition~\ref{prop:F-rows}); the objects are key-averaged transitions and local sums,
not representatives.\\
\bottomrule
\end{tabular}
\end{center}

\subsection{Contributions in three groups}
\begin{enumerate}
\item[(I)] \emph{Differential structure common to all corrections $h$} (Section~\ref{sec:diff}).
Direction-wise counts for the classes $\cZ,\cE,\cG$, the exact uniformity $\Delta(J\circ F)=q$,
the two-block bound with the kernel direction made explicit, a uniform collision-energy bound, and
(Section~\ref{sec:mixing}) the uniform one-word contraction R1 with the exact $\cZ$ block, the $h$-free estimate of the $\cG$
block, the resulting unconditional step count for the averaged alternating arrangement, and the exponent $1/2$ for the
inverse map with $d=1\oplus K_0(x/\rho)$.
\item[(II)] \emph{Constructive correction of local zero sums} (Section~\ref{sec:repair}).
For the inverse map and every non-zero two-to-one affine $d$: a linear certificate on $h$, an explicit correction of Boolean
degree two, the exact pass fraction $1-1/q$ on every linear design space containing it, minimality of
Boolean degree two for this local problem, and the count of cosets on which the local sum is non-zero.
\item[(III)] \emph{Composition guarantees and what they do not prevent} (Section~\ref{sec:comp}).
The averaged multi-word bound for the $(2,2,1,1)$ arrangement, the candidate-existence distinguisher
and its conditioned false-acceptance bound, and the resources $T/D/M$ of the implemented
enumeration.
\end{enumerate}
Group~(III) consists of applications of known principles (key averaging of Markov ciphers
\cite{LMM}, the two-layer MDS count \cite{Hong,Park}, post-composition with a reset key, the zero
sum of all ring elements \cite{KW}, and the RP/RF comparison of \cite{Bernstein}). It is included to
fix the design meaning and the limits of~(I) and~(II), not as an independent novelty claim.

\subsection{What is not claimed}
We do not propose a cipher and do not claim that any construction in this paper is secure.
In particular we do not claim: PRP security of any arrangement; a fixed-key differential bound;
a key schedule; a lower bound on the cost of any attack; that no attack below the codebook exists;
minimality of AND counts of $F$, of $S_t$ or of any arrangement (Proposition~\ref{prop:mc} concerns the correction alone); an advantage in cost over known Feistel or SPN
constructions (for instance a three-round Feistel network with $x^3$ attains the same constant
$4/N$ with fewer AND gates \cite{AO}). The contraction results of Section~\ref{sec:mixing} concern key-averaged difference distributions of the
alternating arrangement; they are not fixed-key statements and do not recommend a number of rounds. Optimality of the exponent
$1/2$ refers to the order of decay in $q$ for the family of Corollary~\ref{cor:half}, not to the constants, to a cipher, or to any other choice of $d$.

\subsection{Related work}
The analysis uses standard tools: branch-number (MDS) arguments \cite{Hong,Park}, Parseval's identity, Newton's identities
for power sums, the Schur test, and the key-averaging view of Markov ciphers \cite{LMM}, under which the expected
differential transition matrix is a Markov matrix whose subdominant spectrum governs convergence to uniformity
(cf.\ \cite{BV}). What is specific here is the object and the explicit, family-uniform bound: for $B=J\circ F$ on $\GR(4,m)$ the
carry-aware difference table restricted to the fluctuation subspace has norm $\tau$ with $\tau^2\le5(3/4)^m-3/q+1/q^2$; in the
sources we read the subdominant eigenvalue is left implicit or assumed, not bounded in this way. The work on sum-free
functions \cite{EHRZ,HZ} concerns sums of a function over $k$-dimensional affine subspaces; here the sums are of the inverse
map over a non-affine image set, their non-vanishing is certified by a linear functional built from reciprocal power sums
over the trace-one hyperplane, and a Boolean-degree-two correction makes the certificate non-zero for every $h$. We did not
find these statements in the works cited; our search was limited, and the techniques themselves are not claimed as new.

% sec_2_setup_v4.tex -- C8 (P8-01, GPT C7-F01): DDT normalized by |X| (N for one word, N^w for w words);
% product rule for word-wise maps stated once. v3: two more displays split (layout only).
\section{Two-digit coordinates, carry, and the PD assumption}\label{sec:setup}

\subsection{Field, ring and coordinates}
Let $m\ge 3$, $q=2^m$, and let $\Fq=\FF[t]/(f)$ with $f$ irreducible of degree $m$. We fix the
polynomial basis $1,t,\dots,t^{m-1}$ once and for all and identify $\Fq$ with $\FF^m$ through it.
We write $\oplus$ for addition in $\Fq$, juxtaposition for the field product, and $a\star u$ for the
coordinate-wise product in the fixed basis (bitwise AND). The operation $\star$ distributes over
$\oplus$ and $u\star u=u$. For $x\in\Fq$, $\wt(x)$ is the number of non-zero coordinates.

Let $R=\GR(4,m)=\mathbb Z_4[t]/(\hat f)$ with $\hat f$ a monic lift of $f$, and $N=|R|=q^2$. Every
$x\in R$ is uniquely $x=\sum_i c_it^i$ with $c_i=a_i+2b_i\in\{0,1,2,3\}$; with $a=(a_i)$,
$b=(b_i)\in\Fq$ and $\sigma(a)=\sum a_it^i$ we have $x=\sigma(a)+2\sigma(b)$ and we write $x=(a,b)$.
The reduction $x\mapsto a$ is the ring homomorphism $R\to\Fq$.

\begin{lemma}[Carry]\label{lem:add}
$(a,b)+(u,v)=(a\oplus u,\ b\oplus v\oplus(a\star u))$ and $-(a,b)=(a,\ a\oplus b)$.
For $y=(y_0,y_1)$, $z=(z_0,z_1)$ and $s=y_0\oplus z_0$,
\[ y-z=\bigl(s,\ y_1\oplus z_1\oplus(z_0\star s)\bigr). \]
\end{lemma}
% SOURCE: c1/proofs/ClassTransition.md:18-22 (Lemma A, A'); negation: c5/a1/notes/UE_P5.md:9
\begin{proof}
Coefficient-wise $a_i+u_i=(a_i\oplus u_i)+2a_iu_i$ in $\mathbb Z_4$, and $2\sigma(x)+2\sigma(y)=2\sigma(x\oplus y)$
because $4=0$. This gives the sum. For the difference, $z+(s,w)=(y_0,\ z_1\oplus w\oplus(z_0\star s))$
by the sum formula, which equals $y$ exactly when $w=y_1\oplus z_1\oplus(z_0\star s)$. The negation
is the case $y=0$.
\end{proof}

In particular a difference with $s=0$ has XOR high digit, while for $s\ne0$ the high digit carries
the correction $z_0\star s$, which takes $2^{\wt(s)}$ values as $z_0$ varies.

\subsection{Polynomial maps in coordinates}
\begin{proposition}[Coordinate form and realizability]\label{prop:realize}
Let $\mathsf F\in R[X]$ and let $\bar{\mathsf F}\in\Fq[X]$ be its reduction. Then
\[
\mathsf F(a,b)=\bigl(\bar{\mathsf F}(a),\ h(a)\oplus d(a)\,b\bigr),\qquad d=\bar{\mathsf F}'\ \text{(formal derivative)},
\]
where $h(a)$ is the high digit of $\mathsf F(\sigma(a))$. Conversely, for arbitrary functions
$p,d,h:\Fq\to\Fq$ there is $\mathsf F\in R[X]$ with $\mathsf F(a,b)=(p(a),h(a)\oplus d(a)b)$.
\end{proposition}
% SOURCE: c1/proofs/ClassTransition.md:25-27 (Lemma T); c3/notes/PD_ROLES_H0_P3.md:19-23 (section 1), :36-38 (section 2.1)
\begin{proof}
For $y,z\in R$, $(y+2z)^n=y^n+2nzy^{n-1}+4(\cdots)$, hence $\mathsf F(y+2z)=\mathsf F(y)+2z\mathsf F'(y)$.
With $y=\sigma(a)$, $z=\sigma(b)$, the term $2\sigma(b)\mathsf F'(\sigma(a))$ equals $2\sigma(\bar{\mathsf F}'(a)b)$, and adding
an element of $2R$ only XORs the high digit. For the converse let $p_{\rm rep},d_{\rm rep}$ be the
interpolating polynomials of degree $<q$ and put $G=p_{\rm rep}+(X^q-X)(d_{\rm rep}+p_{\rm rep}')$. On $\Fq$,
$G=p$, and $G'=p'_{\rm rep}+(qX^{q-1}-1)(d_{\rm rep}+p'_{\rm rep})+(X^q-X)(\cdots)'$ equals
$p'_{\rm rep}+d_{\rm rep}+p'_{\rm rep}=d$ on $\Fq$ (characteristic $2$, $q\equiv0$). Let $\mathsf F_0$ be any
coefficient lift of $G$, with high digit $h_0$. If $\hat H$ lifts an interpolant of $h\oplus h_0$,
then $2\hat H(\sigma(a)+2\sigma(b))=2\sigma(h(a)\oplus h_0(a))$, so $\mathsf F=\mathsf F_0+2\hat H$ has high digit
$h(a)\oplus d(a)b$.
\end{proof}

Thus the quantifier ``over all coefficient lifts'' is the same as ``over all functions $h$'', and
$d$ is not tied to $p$. Let $J(a,b)=(b,a)$; $J$ is a bijection of $R$ which is not additive.
For $p,d,h$ as above put
\begin{gather*}
F=F_{p,d,h}:(a,b)\mapsto (p(a),\,h(a)\oplus d(a)b),\\
B=J\circ F:(a,b)\mapsto(h(a)\oplus d(a)b,\,p(a)).
\end{gather*}

\subsection{The PD assumption}
For $u\in\Fq$ write $D_up(a)=p(a\oplus u)\oplus p(a)$. A map $p$ is APN if for every $u\ne0$ and
every $c$ the equation $D_up(x)=c$ has either no solution or exactly the two solutions
$\{x,x\oplus u\}$.

\begin{assumption}[PD]\label{ass:PD}
\begin{enumerate}
\item[(P1)] $p:\Fq\to\Fq$ is an APN permutation (not necessarily quadratic);
\item[(P2)] $d=\varepsilon\oplus K$ with $K$ $\FF$-linear, $\ker K=\{0,\rho\}$ for some $\rho\ne0$, and $\varepsilon\notin\Img K$;
\item[(P3)] $h:\Fq\to\Fq$ is arbitrary.
\end{enumerate}
We write $\omega=2^{\wt(\rho)}$, where the weight is taken in the fixed basis, and
$S_\rho=\{x\star\rho:x\in\Fq\}$, so $|S_\rho|=\omega$.
\end{assumption}
% SOURCE: c3/notes/PD_ROLES_H0_P3.md:62-73 (theorem PD and its hypotheses); c5/notes/CORE_P5.md:9-13

\begin{lemma}\label{lem:d}
Under (P2): $d(a)\ne0$ for all $a$; $d(a\oplus u)\oplus d(a)=K(u)$; $d(a\oplus\rho)=d(a)$; every
non-empty fibre of $d$ is a pair $\{a,a\oplus\rho\}$; $\Img d=\varepsilon\oplus\Img K$ is an affine
hyperplane of $\Fq$ not containing $0$, so $|\Img d|=q/2$.
\end{lemma}
% SOURCE: c1/proofs/ClassTransition.md:12; c3/notes/PD_ROLES_H0_P3.md:42 (Fib follows from Aff, Ker)
\begin{proof}
$d(a)=0$ would mean $\varepsilon=K(a)\in\Img K$. The other statements follow from linearity of $K$
and $\dim\ker K=1$.
\end{proof}

\begin{remark}
The basis is fixed. Replacing $p$ by $p(\rho x)$ would normalize $\rho=1$, but this changes the map
because $J$ and the carry are defined in the fixed basis; we never do this. For the inverse
$p(x)=x^{q-2}$ with $m$ odd and $d_0(a)=(a^2+a+1)^2$ one has $\rho=1$, $\omega=2$
(Section~\ref{sec:repair}). The inverse map is an APN permutation exactly for odd $m$; when an APN permutation
is needed for all statements at once we therefore take $p(x)=x^{q-2}$ with $m\ge5$ odd. The counting lemmas of
Section~\ref{sec:diff} and Theorem~\ref{thm:R1} use only (P1)--(P3) and do not use the parity of $m$.
\end{remark}

\subsection{Difference distribution and direction classes}
All differences are ring differences. For a finite abelian group $X$ (below $X=R$ or $X=R^w$) and a map
$G:X\to X$ put
\begin{gather*}
\delta_G(\alpha,\beta)=\#\{x\in X: G(x+\alpha)-G(x)=\beta\},\\
T_G=\DP_G=\delta_G/|X|,\qquad \Delta(G)=\max_{\alpha\ne0,\beta}\delta_G(\alpha,\beta).
\end{gather*}
The normalizing factor is $N$ for one word and $N^w$ for $w$ words, so that every row of $T_G$ sums to $1$.
For a word-wise map $G=(G_1,\dots,G_w)$ one has $\delta_G=\prod_i\delta_{G_i}$ and $|R^w|=N^w$, hence
$T_G(\alpha,\beta)=\prod_iT_{G_i}(\alpha_i,\beta_i)$; multi-word tables enter only in this product form
(Proposition~\ref{prop:markov}, Lemmas~\ref{lem:reset}, \ref{lem:D1}, Section~\ref{sec:mixing}). If $G$ is a permutation, $T_G$ is doubly stochastic, $T_G(0,\cdot)=e_0$ and $T_G(\alpha,0)=0$ for $\alpha\ne0$.
A non-zero difference $\alpha=(u,v)$ belongs to the class
\[
\cZ:\ u=0\ (q-1\text{ elements}),\qquad \cE:\ u=\rho\ (q),\qquad \cG:\ u\notin\{0,\rho\}\ (q(q-2)).
\]

\subsection{Keys, local blocks and averaged probabilities}
Keys are added in $R$. For $t\ge1$ the local block with $t$ key words is
\[
S_t(x;k_1,\dots,k_t)=F\bigl(B(\cdots B(B(x)+k_1)\cdots)+k_t\bigr),
\]
so $S_t$ uses $t+1$ evaluations of $F$ and $t$ of $J$. In the averaging model the keys, and a
whitening key $k_0$ in front, are independent and uniform.

\begin{proposition}[Key averaging; standard]\label{prop:markov}
For $\alpha,\gamma\in R$,
\[\E_{k_0,\dots,k_t}\DP_{x\mapsto S_t(x+k_0;k)}(\alpha,\gamma)=(T_B^{\,t}T_F)(\alpha,\gamma).\]
\end{proposition}
% SOURCE: c1/notes/MODEL.md:22,29 (model and key-averaged DP); c2/notes/AVERAGED_KERNEL_MDS.md:20-25. Known principle [LMM].
\begin{proof}
This is the Markov-cipher principle \cite{LMM}: for fixed inputs $x,x+\alpha$, the value $x+k_0$
is uniform and the difference is $\alpha$; after each $B$ the addition of an independent uniform key
makes the next input uniform and independent of the difference already produced, so the
difference evolves by the transition matrix $T_B$, and finally by $T_F$.
\end{proof}

We write $\EDP_t=\max_{\alpha\ne0,\gamma}(T_B^tT_F)(\alpha,\gamma)$.

\subsection{The \texorpdfstring{$(2,2,1,1)$}{(2,2,1,1)} arrangement}\label{ss:arr}
Let $w\ge2$, let $M_1,M_2,M_3$ be $R$-linear automorphisms of $R^w$, and let local layers act
word-wise. The arrangement studied in Section~\ref{sec:comp} is
\[
E(x)=S_1^{(4)}\Bigl(M_3\,S_1^{(3)}\bigl(M_2\,S_2^{(2)}(M_1\,S_2^{(1)}(x+k_0)+L_1)+L_2\bigr)+L_3\Bigr),
\]
with state names $X_1=S_2^{(1)}(x+k_0)$, $V_1=M_1X_1+L_1$, $X_2=S_2^{(2)}(V_1)$, $V_2=M_2X_2+L_2$,
$X_3=S_1^{(3)}(V_2)$, $V_3=M_3X_3+L_3$, $Y=S_1^{(4)}(V_3)$. It has three mixing layers and two
reset keys $L_2,L_3$ in front of the one-key layers. The key words are
$k_0,L_1,L_2,L_3,K_3,K_4$ ($w$ words each) and $K_1,K_2$ ($2w$ words each), $10w$ words in total,
all independent and uniform in the averaging model. These $10w$ words are the randomness of the
model, not a security level; no key schedule is specified. The arrangement and the number of mixing
layers are fixed throughout; we do not add layers. We assume that every entry and every minor of
$M_1$ is a unit (MDS) and that $M_2,M_3$ are invertible.
% SOURCE: c4/notes/CORE_P4.md:49-55; c5/spec/CANDIDATE_P5_SPLIT_TAIL.json (structure, keys, mixers, claims_not_made)

% sec_3_differential_v3.tex -- v3: one more display and the table of prop:table (layout only). Group I.
\section{Differential structure common to all corrections}\label{sec:diff}

Throughout this section Assumption~\ref{ass:PD} holds, $F=F_{p,d,h}$ and $B=J\circ F$.
No statement depends on $h$.

\subsection{Difference formulas}
\begin{lemma}\label{lem:diffF}
For $x=(a,b)$ and $\alpha=(u,v)$,
\begin{multline*}
F(x+\alpha)-F(x)=\Bigl(D_up(a),\\ h(a\oplus u)\oplus h(a)\oplus K(u)b\oplus d(a\oplus u)(v\oplus a\star u)\oplus\bigl(p(a)\star D_up(a)\bigr)\Bigr).
\end{multline*}
\end{lemma}
% SOURCE: c1/proofs/ClassTransition.md:31-35 (formula (0.1)); PD version c3/notes/PD_ROLES_H0_P3.md:47-51
\begin{proof}
By Lemma~\ref{lem:add}, $x+\alpha=(a\oplus u,\ b\oplus v\oplus a\star u)$. Apply the difference formula
of Lemma~\ref{lem:add} with $z_0=p(a)$, $s=D_up(a)$, and use $d(a\oplus u)b\oplus d(a)b=K(u)b$.
\end{proof}

\begin{lemma}\label{lem:diffB}
For $x=(a,b)$ and $\alpha=(u,v)$ put $Q=h(a)\oplus d(a)b$. Then $B(x+\alpha)-B(x)=(s,t)$ with
\[
s=h(a\oplus u)\oplus h(a)\oplus K(u)\,b\oplus d(a\oplus u)(v\oplus a\star u),\qquad t=D_up(a)\oplus(Q\star s).
\]
\end{lemma}
% SOURCE: c2/notes/EPSILON_THEOREM.md:13-14; c1/proofs/ClassTransition.md:45-46,65; c5/a1/notes/UE_P5.md:14-18
\begin{proof}
$B(x)=(Q,p(a))$ and $B(x+\alpha)=(Q',p(a\oplus u))$ with
$Q'=h(a\oplus u)\oplus d(a\oplus u)(b\oplus v\oplus a\star u)$. Lemma~\ref{lem:add} with $z_0=Q$ and
$s=Q'\oplus Q$ gives the claim.
\end{proof}

The output class of a difference of $B$ is read off its low digit $s$: $s=0$, $s=\rho$, or
$s\notin\{0,\rho\}$. Since $B$ is a permutation (by (P1) and $d\ne0$), $\alpha\ne0$ gives $(s,t)\ne0$.

\begin{proposition}[Rows of $T_F$]\label{prop:F-rows}
Let $\alpha=(u,v)\ne0$.
\begin{enumerate}
\item If $\alpha\in\cG$ then $\delta_F(\alpha,\beta)\le2$ for all $\beta$.
\item If $\alpha\in\cZ$ then $F(x+\alpha)-F(x)=(0,d(a)v)$; hence $\delta_F(\alpha,(0,w))=2q$ if
$w/v\in\Img d$ and $\delta_F(\alpha,\beta)=0$ otherwise.
\item If $\alpha\in\cE$ then the high digit of $F(x+\alpha)-F(x)$ does not depend on $b$ and, for fixed $a$, is a bijective function of $v$; in particular $\delta_F(\alpha,\beta)\le2q$.
\end{enumerate}
\end{proposition}
% SOURCE: c1/proofs/ClassTransition.md:37-41 (Lemma 0 (i)-(iii)); c3/notes/PD_ROLES_H0_P3.md:49-51
\begin{proof}
(1) $K(u)\ne0$, so for each $a$ at most one $b$ gives a prescribed high digit, and $D_up(a)=\bar\beta$
has at most two solutions by (P1). (2) Put $u=0$ in Lemma~\ref{lem:diffF}; $d$ is two-to-one onto
$\Img d$. (3) $K(\rho)=0$ removes $b$, and the coefficient of $v$ is $d(a\oplus\rho)=d(a)\ne0$; at most two $a$
solve $D_\rho p(a)=\bar\beta$, each with $q$ values of $b$.
\end{proof}

\subsection{Exact differential uniformity}
\begin{theorem}\label{thm:eps}
$\Delta(B)=q$; that is, $\max_{\alpha\ne0,\beta}T_B(\alpha,\beta)=1/q$. More precisely, for every
$\beta$: $\delta_B(\alpha,\beta)\le q$ if $\alpha\in\cG$; $\delta_B(\alpha,\beta)\le 2q/2^{\wt(s)}\le q$ if
$\alpha\in\cZ$ and $\beta=(s,t)$; and $\delta_B(\alpha,\beta)\le q$ if $\alpha\in\cE$. The value $q$ is attained.
\end{theorem}
% SOURCE: c2/notes/EPSILON_THEOREM.md:9-22 (theorem E, proof cases 1-4); generality: c3/notes/PD_ROLES_H0_P3.md:57,70
\begin{proof}
Use Lemma~\ref{lem:diffB} with $\beta=(s,t)$.

\emph{$\cG$.} $K(u)\ne0$, so for fixed $a$ the value $s$ determines $b$. Hence at most one $b$ per $a$, at most $q$ in total.

\emph{$\cZ$.} Here $s=d(a)v\ne0$, so at most two $a$ (fibres of $d$). For fixed $a$, $Q$ is a
bijective function of $b$ and $t=Q\star s$; the map $Q\mapsto Q\star s$ takes $2^{\wt(s)}$ values, each
$q/2^{\wt(s)}$ times. Total $\le 2q/2^{\wt(s)}\le q$ as $s\ne0$.

\emph{$\cE$.} $K(\rho)=0$, so $s=s(a)$ does not depend on $b$, and from $d(a\oplus\rho)=d(a)$ and
$(a\oplus\rho)\star\rho=(a\star\rho)\oplus\rho$,
\begin{equation}\label{eq:C3}
s(a\oplus\rho)=s(a)\oplus d(a)\rho\ne s(a).
\end{equation}
Let $S_s=\{z\star s:z\in\Fq\}$. The condition is $s(a)=s$ and $D_\rho p(a)\in t\oplus S_s$. By (P1) every
non-empty fibre of $D_\rho p$ is a pair $\{a,a\oplus\rho\}$, so the admissible $a$ lie in at most $|S_s|$
pairs; by~\eqref{eq:C3} each pair contains at most one $a$ with $s(a)=s$; for that $a$ exactly
$q/|S_s|$ values of $b$ give the required $Q\star s$. Total $\le q$.

\emph{Attainment.} Let $\eta$ be a basis vector, $d_*\in\Img d$, $v=\eta/d_*$, $\alpha=(0,v)$, $\beta=(\eta,0)$.
Then $s=\eta$ exactly for the two $a$ with $d(a)=d_*$, and for each such $a$ exactly $q/2$ values of $b$
give $Q\star\eta=0$.
\end{proof}

The proof uses neither $\wt(\rho)$ nor the permutation property of $p$ in the cases $\cZ,\cE$.

\subsection{Class transitions}
\begin{proposition}\label{prop:table}
For $\beta$ in the indicated class, the following column sums (over all rows $\alpha$ of the indicated
class) and entry bounds hold for $T_B$:
\begin{center}
\begin{tabular}{llp{0.12\textwidth}p{0.36\textwidth}}
\toprule
rows $\to$ target & quantity & bound & ingredient\\
\midrule
$\cZ\to\cZ$ & column sum & $0$ & $s=d(a)v\ne0$\\
$\cZ\to\cE$ & column sum & $1/\omega$ & fibres of $d$, $Q$ uniform in $b$\\
$\cE\to\cZ$ & column sum & $2/q$ & $v\mapsto s_v(a)$ bijective, APN in direction $\rho$\\
$\cE\to\cE$ & column sum & $2/q$ & as above, $D_\rho p(a)\in t\oplus S_\rho$\\
$\cG\to\cZ$ & entry & $2/N$ & $b$ unique per $a$, APN\\
$\cG\to\cE$ & entry & $2\omega/N$ & $b$ unique per $a$, $D_up(a)\in t\oplus S_\rho$\\
\bottomrule
\end{tabular}
\end{center}
For the final layer $F$: if $\gamma$ has low digit $0$, only rows in $\cZ$ contribute and their column
sum is at most $1$; if $\gamma$ has non-zero low digit, rows in $\cZ$ contribute $0$, rows in $\cE$ have
column sum at most $2/q$, and rows in $\cG$ have entries at most $2/N$.
\end{proposition}
% SOURCE: c1/proofs/ClassTransition.md:63-78 (section 4 table and final layer); PD generality c3/notes/PD_ROLES_H0_P3.md:55
\begin{proof}
All counts use Lemma~\ref{lem:diffB} and divide by $N=q^2$.
$\cZ\to\cZ$: $s=d(a)v\ne0$.
$\cZ\to\cE$, target $(\rho,t)$: for each $v$ with $\rho/v\in\Img d$ (there are $q/2$ such $v$) exactly two
$a$ have $s=\rho$; then $t=Q\star\rho$ and $Q$ is uniform in $b$, so $q/\omega$ values of $b$ (if $t\in S_\rho$).
The column sum is at most $(q/2)\cdot2\cdot(q/\omega)/N=1/\omega$.
$\cE\to\cZ$, target $(0,t)$: for fixed $a$, $v\mapsto s_v(a)=h(a\oplus\rho)\oplus h(a)\oplus d(a)(v\oplus a\star\rho)$ is
bijective, so one $v$ per $a$ gives $s=0$; then $t=D_\rho p(a)$ has at most two solutions $a$, each with
$q$ values of $b$: column sum $\le 2q/N$.
$\cE\to\cE$, target $(\rho,t)$: one $v$ per $a$; $D_\rho p(a)\in t\oplus S_\rho$ for at most $2\omega$ values of $a$,
each with $q/\omega$ values of $b$: column sum $\le2q/N$.
$\cG\to\cZ$: $s=0$ fixes $b$ for each $a$, and $t=D_up(a)$ has at most two solutions.
$\cG\to\cE$: $s=\rho$ fixes $b$, and $D_up(a)\in t\oplus S_\rho$ has at most $2\omega$ solutions.
Final layer: by Lemma~\ref{lem:diffF} the low digit of an $F$-difference is $D_up(a)$, which is non-zero
for $u\ne0$ because $p$ is a permutation; $T_F$ is doubly stochastic, so any partial column sum is at most
$1$. For rows in $\cE$ apply Proposition~\ref{prop:F-rows}(3): at most two $a$, one $v$ per $a$, $q$ values
of $b$. For rows in $\cG$ apply Proposition~\ref{prop:F-rows}(1).
\end{proof}

\subsection{A sup-norm recursion and the two-block bound}
Fix $\alpha\ne0$ and let $\pi_t=e_\alpha T_B^{\,t}$ be the distribution of the difference after $t$
key-separated applications of $B$ (Proposition~\ref{prop:markov}). Put
$\sigma^{\cZ}_t=\max_{\beta\in\cZ}\pi_t(\beta)$ and $\sigma^{\cE}_t=\max_{\beta\in\cE}\pi_t(\beta)$.

\begin{theorem}[Sup-norm recursion]\label{thm:S}
For every $t\ge0$,
\begin{gather*}
\sigma^{\cZ}_{t+1}\le\frac2N+\frac2q\,\sigma^{\cE}_t,\qquad
\sigma^{\cE}_{t+1}\le\frac{2\omega}N+\frac1\omega\,\sigma^{\cZ}_t+\frac2q\,\sigma^{\cE}_t,\\
\max_\gamma(\pi_tT_F)(\gamma)\le\max\Bigl(\sigma^{\cZ}_t,\ \frac2N+\frac2q\,\sigma^{\cE}_t\Bigr).
\end{gather*}
\end{theorem}
% SOURCE: c1/proofs/AllDirectionBound.md:26-31 (theorem S); constants c1/proofs/ClassTransition.md:63-78
\begin{proof}
$\pi_{t+1}(\beta)=\sum_{\alpha'}\pi_t(\alpha')T_B(\alpha',\beta)$ with $\pi_t(0)=0$. Split $\alpha'$ by class. The rows in
$\cG$ contribute at most (largest entry)$\cdot\pi_t(\cG)\le$ largest entry; the rows in $\cZ$ (resp.\ $\cE$)
contribute at most $\sigma^{\cZ}_t$ (resp.\ $\sigma^{\cE}_t$) times the column sum over that class. Insert
Proposition~\ref{prop:table}. The last inequality is the same argument with the final-layer part of
Proposition~\ref{prop:table}.
\end{proof}

\begin{theorem}[Two-block bound, kernel direction explicit]\label{thm:W2}
Let $m\ge5$ be odd. For every $h$ and every $\alpha\ne0$, $\gamma$:
\begin{enumerate}
\item if $\wt(\rho)=1$ (so $\omega=2$), then $(T_B^2T_F)(\alpha,\gamma)\le 4/N$;
\item for general $\rho$, $(T_B^2T_F)(\alpha,\gamma)\le(6+6/q)/N\le 99/(16N)$.
\end{enumerate}
By Proposition~\ref{prop:markov}, $\EDP_2$ of the local block $S_2$ obeys the same bounds.
\end{theorem}
% SOURCE: c3/notes/MAINLINE_P3.md:6-10 (lemma W, warm start, theorem W2), :14-16 (general rho); scope correction c5/notes/CORE_P5.md:5-13
\begin{proof}
Put $x=1/q$, so $x^2=1/N$. By Theorem~\ref{thm:eps}, $\pi_1(\beta)=T_B(\alpha,\beta)\le x$ for all $\beta$, hence
$\sigma^{\cZ}_1,\sigma^{\cE}_1\le x$. Theorem~\ref{thm:S} uses of $\pi_t$ only that it is a probability vector and
the two maxima, so it may be started at $t=1$:
$\sigma^{\cZ}_2\le 2x^2+2x\cdot x=4x^2$ and $\sigma^{\cE}_2\le 2\omega x^2+x/\omega+2x^2$. Therefore
$(T_B^2T_F)(\alpha,\gamma)\le\max(4x^2,\ 2x^2+2x\sigma^{\cE}_2)$ and
\[
N\bigl(2x^2+2x\sigma^{\cE}_2\bigr)=2+\frac2\omega+\frac{4(\omega+1)}q .
\]
For $\omega=2$ this is $3+12/q\le4$ since $q\ge12$. For general $2\le\omega\le q$,
$(6+6/q)-\bigl(2+2/\omega+4(\omega+1)/q\bigr)=(4\omega-2)(q-\omega)/(q\omega)\ge0$, and $4\le6+6/q$.
Finally $6+6/q\le 99/16$ for $q\ge32$.
\end{proof}

\begin{remark}
The constant $4$ in (1) uses $\omega=2$, i.e.\ weight one of $\rho$ in the fixed basis; it is not
obtained for general $\rho$ by a change of basis (Section~\ref{sec:setup}). All tables used in
Sections~\ref{sec:repair}--\ref{sec:comp} have $\rho=1$.
% SOURCE: c5/notes/CORE_P5.md:13-15 (rho of the actual tables)
\end{remark}

\subsection{A uniform collision-energy bound}
Define the collision energy of $B$ as $\mathcal E_B=\sum_{\alpha\in R}\sum_{\beta}T_B(\alpha,\beta)^2$ (the row
$\alpha=0$ included).

\begin{theorem}[Uniform energy]\label{thm:UE}
Assume (P2), (P3) and that $p$ is APN (the permutation property is not used). Then for every $m$,
every $h$ and every $\rho$,
\[
\mathcal E_B\le 6+4(3/4)^m,\qquad\text{and}\qquad \mathcal E_B\le 1779/256<7\ \text{ if } m\ge5 .
\]
\end{theorem}
% SOURCE: c5/a1/notes/UE_P5.md:69-75 (theorem UE); proof :14-68
The proof occupies the rest of this subsection. Write $\mathcal E_{\cZ},\mathcal E_{\cE},\mathcal E_{\cG}$ for the
contributions of the rows in each class; the zero row contributes $1$.

\begin{lemma}[$\cZ$ rows, exact]\label{lem:UEZ}
$\mathcal E_{\cZ}=2(3/4)^m-2/q$.
\end{lemma}
% SOURCE: c5/a1/notes/UE_P5.md:63-65
\begin{proof}
For $u=0$, Lemma~\ref{lem:diffB} gives $s=d(a)v$ and $t=Q\star s$. Hence
$\delta_B((0,v),(s,t))=2q/2^{\wt(s)}$ if $s\in v\,\Img d$ and $t\subseteq s$ (bitwise), and $0$ otherwise. The
row energy is $(4/q^2)\sum_{s\in v\Img d}2^{-\wt(s)}$. Each $s\ne0$ lies in $v\Img d$ for exactly $|\Img d|=q/2$
values of $v\ne0$. Thus $\mathcal E_{\cZ}=(4/q^2)(q/2)\sum_{s\ne0}2^{-\wt(s)}=(2/q)((3/2)^m-1)$.
\end{proof}

\begin{lemma}[$\cE$ rows]\label{lem:UEE}
$\mathcal E_{\cE}\le1$.
\end{lemma}
% SOURCE: c5/a1/notes/UE_P5.md:66
\begin{proof}
By the $\cE$ case of Theorem~\ref{thm:eps}, which uses only APN and (P2), every entry of an $\cE$ row is at
most $1/q$; rows sum to $1$, so each row energy is at most $1/q$, and there are $q$ rows.
\end{proof}

For the $\cG$ rows fix $u\notin\{0,\rho\}$ and put
\[
\begin{gathered}
\lambda(a)=d(a)K(u)^{-1},\qquad w(a)=K(u)\lambda(a)\lambda(a\oplus u),\\
H_u(a)=\lambda(a\oplus u)h(a)\oplus\lambda(a)h(a\oplus u).
\end{gathered}
\]
\begin{lemma}\label{lem:Grow}
For $\alpha=(u,v)$ and each $(a,s)$ there is exactly one $b$ with low difference $s$, and then
\begin{gather*}
\delta_B((u,v),(s,t))=\#\{a: t_s(a)=t\},\qquad t_s(a)=D_up(a)\oplus\bigl(s\star z_s(a)\bigr),\\
z_s(a)=H_u(a)\oplus\lambda(a)s\oplus w(a)(v\oplus a\star u).
\end{gather*}
Moreover:
\textup{(S1)} $\lambda(a\oplus\rho)=\lambda(a)$, $\lambda(a\oplus u)=\lambda(a)\oplus1$, $w(a\oplus\rho)=w(a\oplus u)=w(a)\ne0$;
\textup{(S2)} $w(a')=w(a)$ iff $a'\in a\oplus\{0,\rho,u,u\oplus\rho\}$;
\textup{(S3)} $z_s(a)\oplus z_s(a\oplus u)=s\oplus w(a)u$; $z_s(a)\oplus z_s(a\oplus\rho)=Z_\rho(a)$ and
$z_s(a)\oplus z_s(a\oplus\rho\oplus u)=s\oplus Z_{\rho u}(a)$, where $Z_\rho(a),Z_{\rho u}(a)$ depend neither on $s$ nor on $v$.
\end{lemma}
% SOURCE: c5/a1/notes/UE_P5.md:22-32 (facts S1-S3)
\begin{proof}
Solve $s$ in Lemma~\ref{lem:diffB} for $b$ (possible as $K(u)\ne0$) and substitute in $Q$; use
$\lambda(a\oplus u)=(d(a)\oplus K(u))K(u)^{-1}=\lambda(a)\oplus1$, so $h(a)(1\oplus\lambda(a))=\lambda(a\oplus u)h(a)$ and
$\lambda(a)d(a\oplus u)=w(a)$. (S1) follows from Lemma~\ref{lem:d}; $w(a)\ne0$ because $d$ has no zeros.
(S2): $w=K(u)g(\lambda)$ with $g(x)=x^2+x$, which is two-to-one with $g(x)=g(x\oplus1)$; so $w(a')=w(a)$ iff
$d(a')\in\{d(a),d(a\oplus u)\}$ iff $a'\oplus a\in\ker K$ or $a'\oplus a\oplus u\in\ker K$.
(S3): $H_u(a\oplus u)=H_u(a)$ and $(a\star u)\oplus((a\oplus u)\star u)=u$; in the other two differences
$\lambda$ and $w$ take the same value at both points, so the $\lambda(a)s$ terms and the $w(a)v$ terms cancel
(resp.\ leave exactly $s$).
\end{proof}

\begin{lemma}[Energy identity]\label{lem:Eid}
With $A_u=\sum_a 2^{\wt(w(a)u)}$,
$\mathrm{Merge}_u=\#\{(s,a,a'): a'\in\{a\oplus\rho,a\oplus\rho\oplus u\},\ t_s(a)=t_s(a')\}$ (independent of $v$) and
$\mathrm{Cross}_u=\sum 3^{\,m-\wt(r)}$ over ordered pairs $(a,a')$ with $w(a)\ne w(a')$, $r=D_up(a)\oplus D_up(a')$,
\[
\sum_v\sum_{s,t}\delta_B((u,v),(s,t))^2=q^3+q\,A_u+\mathrm{Cross}_u+q\,\mathrm{Merge}_u .
\]
\end{lemma}
% SOURCE: c5/a1/notes/UE_P5.md:34-45 (identity E, definition of Merge)
\begin{proof}
For each $s$ the row has exactly $q$ points with low difference $s$ (one per $a$), so
$\sum_t\delta_B^2=\#\{(a,a'):t_s(a)=t_s(a')\}$. Split the ordered pairs:
(i) $a'=a$ gives $q^3$ after summing over $v,s,a$.
(ii) $a'=a\oplus u$: $D_up(a\oplus u)=D_up(a)$ and the condition is $s\star(s\oplus w(a)u)=0$, i.e.\ $s\subseteq w(a)u$,
which holds for $2^{\wt(w(a)u)}$ values of $s$ and every $v$.
(iii) $a'\in\{a\oplus\rho,a\oplus\rho\oplus u\}$: by (S3) the condition does not depend on $v$.
(iv) otherwise $w(a)\ne w(a')$ by (S2), and $z_s(a)\oplus z_s(a')=c\oplus(w(a)\oplus w(a'))v$ with $c$ independent
of $v$; the condition $r=s\star(c\oplus(w\oplus w')v)$ has $2^{m-\wt(s)}$ solutions $v$ if $r\subseteq s$ and none
otherwise, and $\sum_{s\supseteq r}2^{m-\wt(s)}=3^{m-\wt(r)}$.
\end{proof}

\begin{lemma}\label{lem:UEbounds}
$\mathrm{Merge}_u\le q^2$, $A_u\le4(3^m-1)$, $\mathrm{Cross}_u\le 2q(q^2-3^m)$.
\end{lemma}
% SOURCE: c5/a1/notes/UE_P5.md:47-58 (lemmas M, A, X)
\begin{proof}
\emph{Merge.} Let $r_a=D_up(a)\oplus D_up(a\oplus\rho)$. Since $a\oplus\rho\notin\{a,a\oplus u\}$, APN gives
$r_a\ne0$, and $D_up(a\oplus\rho\oplus u)=D_up(a\oplus\rho)$ gives the same $r_a$ for the second partner. For
$a'=a\oplus\rho$ the condition is $r_a=s\star Z_\rho(a)$, with $2^{m-\wt(Z_\rho)}\le q/2$ solutions (as $\wt(Z_\rho)\ge\wt(r_a)\ge1$
whenever solutions exist). For $a'=a\oplus\rho\oplus u$ the condition is $r_a=s\star\overline{Z_{\rho u}}$ (bitwise
complement), with $2^{\wt(Z_{\rho u})}\le q/2$ solutions. Hence $\mathrm{Merge}_u\le q\cdot2\cdot q/2$.
\emph{$A_u$.} $w(a)u\ne0$ and, by (S2), each value of $a\mapsto w(a)u$ has at most $4$ preimages; so
$A_u\le4\sum_{z\ne0}2^{\wt(z)}=4(3^m-1)$.
\emph{Cross.} In case (iv), $r\ne0$ (otherwise $a'\in\{a,a\oplus u\}$ by APN). For fixed $a$ and $r\ne0$ there are
at most two $a'$; so $\mathrm{Cross}_u\le2q\sum_{r\ne0}3^{m-\wt(r)}=2q(4^m-3^m)$.
\end{proof}

\begin{proof}[Proof of Theorem~\ref{thm:UE}]
By Lemmas~\ref{lem:Eid} and~\ref{lem:UEbounds}, each $u\notin\{0,\rho\}$ contributes at most
$(4q^3+2q\,3^m-4q)/q^4$, and there are $q-2$ such $u$, so
$\mathcal E_{\cG}\le(1-2/q)(4+2(3/4)^m-4/q^2)\le4+2(3/4)^m$. Adding the zero row, Lemma~\ref{lem:UEZ}
and Lemma~\ref{lem:UEE} gives $\mathcal E_B\le 1+2(3/4)^m+1+4+2(3/4)^m$. For $m\ge5$,
$6+4\cdot243/1024=1779/256$.
\end{proof}

\begin{remark}
Theorem~\ref{thm:UE} gives a constant independent of $m,p,d,h$; it does not by itself give a mixing
statement or any security statement. Its use in the conditional mixing analysis is recorded in
Section~\ref{sec:mixing}.
\end{remark}

% sec_4_repair_v6.tex -- C9 paper lane (P9-09/P9-10): two internal references in the text removed
% ("the C6 notes", "of C6"); mathematics unchanged from sec_4_repair_v5.tex.
% sec_4_repair_v5.tex -- C8: P8-05 (quantifier of non-necessity, GPT C7-F03; Example ex:quant),
% P8-06 (Proposition prop:mc, multiplicative complexity m-1 of g_d alone).
% sec_4_repair_v4.tex (v4: one display split, layout only) -- Group II, written by the main author (C7, P7-02..P7-05, P7-14).
% Replaces the frame sec_4_repair_v2.tex. The d0 statements of sec_4_d0_notes_v2.tex appear here as
% specializations (Corollary cor:d0); labels dz:prop:link and dz:prop:X21 are kept for Section 8.
% Proof sources: c7/notes/resp/P7-02.md .. P7-05.md (paper), GPT NEW_DERIVATIONS_JA.md A (audit proposal,
% re-proved independently), C6 c6/h/notes/H_CERT_P6.md (d0 case). Finite calibration: c7/gd/.
\section{Local zero sums and their certified correction}\label{sec:repair}

Throughout this section $p(x)=x^{q-2}$ (with $0\mapsto0$), $m\ge3$, and $d=\varepsilon\oplus K$ is any
$\FF$-affine map with $\ker K=\{0,\rho\}$, $\rho\ne0$, $\varepsilon\notin\Img K$ (condition (P2) of
Assumption~\ref{ass:PD}). Oddness of $m$ is not used in Subsections~\ref{ss:norm}--\ref{ss:deg}; it enters
only when the inverse map is required to be APN, i.e.\ when the results are combined with
Section~\ref{sec:diff} (Remark~\ref{rem:same}).

\subsection{The local sum}
For an $\FF$-linear map $L:\Fq\to\Fq$, $\mu\in\Fq$ and $c\in\Fq$ put
\[
\Psi_{L,\mu}(c)=\sum_{a\in\Fq}p\bigl(h(a)\oplus d(a)(c\oplus L(a))\oplus\mu\bigr).
\]
\begin{proposition}[Identification with the local coset sum of $S_2$]\label{dz:prop:link}
Let $F=F_{p,d,h}$ with $(p,d,h)$ as above, $S_2(x;K_1,K_2)=F(B(B(x)+K_1)+K_2)$, $K_1=(\kappa,\lambda)$,
$K_2=(\mu,\nu)$, and let $\mathrm{lo}\,s(a)$ be the exclusive or of the low digits of $S_2((a,b);K_1,K_2)$ over
$b\in\Fq$ (equivalently the low digit of the ring sum over the coset $(a,\ast)$). Then, with
$L(u)=u\star\kappa$ and $c=p(a)\oplus\lambda\oplus\kappa$, $\mathrm{lo}\,s(a)=\Psi_{L,\mu}(c)$. The key word $\nu$
does not occur, $\lambda$ only translates $c$, and $a\mapsto c$ is bijective; hence
$\#\{a:\mathrm{lo}\,s(a)\ne0\}=\#\{c:\Psi_{L,\mu}(c)\ne0\}$.
\end{proposition}
\begin{proof}
$B(a,b)=(h(a)\oplus d(a)b,\ p(a))$. Adding $K_1$ (Lemma~\ref{lem:add}) gives low digit
$u=h(a)\oplus d(a)b\oplus\kappa$ and high digit $p(a)\oplus\lambda\oplus((u\oplus\kappa)\star\kappa)=c\oplus L(u)$.
Since $d(a)\ne0$ (Lemma~\ref{lem:d}), $b\mapsto u$ is a bijection of $\Fq$, so the exclusive or over $b$ is
an exclusive or over $u$. The low digit of $B(u,c\oplus L(u))+K_2$ is $h(u)\oplus d(u)(c\oplus L(u))\oplus\mu$
($\nu$ enters only the high digit), and the low digit of the final $F$ is $p$ of it. Summing over $u$
gives $\Psi_{L,\mu}(c)$. Finally $p$ is a permutation, so $a\mapsto c$ is bijective.
\end{proof}
% SOURCE: c6/notes/resp/P6-05_link.md:11-27 (d0); the derivation uses only d(a)!=0 and p bijective,
%         so it holds for every d satisfying (P2) (c7/notes/resp/P7-03.md).
Thus $\Psi_{L,\mu}\not\equiv0$ means that the local zero sum fails on at least one coset, and a lower bound on
$\#\{c:\Psi\ne0\}$ is a lower bound on the number of cosets on which it fails. The converse direction
(non-vanishing on every coset) is not claimed.

\subsection{Normalization of $d$}\label{ss:norm}
\begin{lemma}\label{lem:norm}
There is a unique $\eta\in\Fq^\ast$ with $\Img d=\{z:\Tr(\eta z)=1\}$. Put $\beta=\eta^{-1}$ and
$r(a)=(\eta d(a))^{q/2}$, the square root of $\eta d(a)$. Then $r$ is $\FF$-affine, $d=\beta r^2$,
$r(a\oplus\rho)=r(a)$, and $r$ maps $\Fq$ two-to-one onto $H=\{x:\Tr(x)=1\}$; the fibre of $r$ over each
$x\in H$ is a pair $\{a,a\oplus\rho\}$. Moreover $0\notin H$.
\end{lemma}
\begin{proof}
$\Img K$ is an $\FF$-hyperplane, so $\Img K=\{z:\Tr(\eta z)=0\}$ for a unique $\eta\ne0$ (non-degeneracy of
the trace form). $\Img d=\varepsilon\oplus\Img K$ is the other coset because $\varepsilon\notin\Img K$, i.e.\
$\Tr(\eta z)=1$ on $\Img d$; uniqueness of $\eta$ follows since a coset of a hyperplane determines the
hyperplane and its defining functional. Squaring is $\FF$-linear and bijective, so $r$ is affine and
$r^2=\eta d$. $\Tr(r(a))=\Tr(r(a)^2)=\Tr(\eta d(a))=1$, so $r(\Fq)\subseteq H$; $d$ is two-to-one with fibres
$\{a,a\oplus\rho\}$ (Lemma~\ref{lem:d}), hence so is $r$, and $|r(\Fq)|=q/2=|H|$. $\Tr(0)=0$ gives $0\notin H$.
\end{proof}
For $d=d_0=(a^2+a+1)^2$ and odd $m$ one has $\eta=\beta=1$, $r(a)=a^2+a+1$, $\rho=1$ (the setting of
\cite{FukuiFFA}). Lemma~\ref{lem:norm} uses no change of basis: $r$ is defined from $d$
and the carry $\star$ is untouched.

\begin{lemma}[Power sums over $H$; known]\label{lem:P}
For $m\ge3$, $p_j:=\sum_{x\in H}x^{-j}$ satisfies $(p_1,p_2,p_3,p_4,p_5)=(1,1,0,1,0)$.
\end{lemma}
\begin{proof}
$Q(X)=1+\sum_{i<m}X^{2^i}$ has $Q'=1$, $Q(0)=1$, $\deg Q=q/2$ and root set $H$, so
$Q(X)=\prod_{x\in H}(1+X/x)$ and the elementary symmetric functions of $\{1/x\}$ are $e_1=e_2=e_4=1$,
$e_3=e_5=0$ ($e_5=0$ also for $m=3$, where $\deg Q=4$). Newton's identities in characteristic two give
$p_1=1$, $p_2=p_1^2=1$, $p_3=e_1p_2+e_2p_1+e_3=0$, $p_4=p_2^2=1$, $p_5=e_1p_4+e_2p_3+e_3p_2+e_4p_1+e_5=0$.
\end{proof}
% Known tool: reciprocal power sums over affine subspaces, [EHRZ] Lemma 4.5 (and Carlitz); not claimed as new.

\subsection{The linear certificate}
For $x=r(a)\in H$ put $\Delta_h(x)=h(a)\oplus h(a\oplus\rho)$ (independent of the choice in the fibre) and
\[
A_d(h)=\sum_{x\in H}x^{-3}\Delta_h(x),\qquad B_d(h)=\sum_{x\in H}x^{-5}\Delta_h(x)^2,
\]
\[
\ell_d(h)=A_d(h)+B_d(h)^{q/2},\qquad \tau=L(\rho).
\]
Let $s=q/2-1$, $z(a)=h(a)\oplus d(a)L(a)\oplus\mu$ and $\Phi(C)=\sum_a(z(a)+d(a)C)^s\in\Fq[C]$.

\begin{theorem}[Top coefficients]\label{thm:top}
For every $h$, $L$, $\mu$: $\Psi_{L,\mu}(c)=\Phi(c)^2$ for all $c\in\Fq$, and
\begin{gather*}
[C^{s}]\Phi=0,\qquad [C^{s-1}]\Phi=\beta^{s-1}\bigl(A_d(h)+\beta\tau\bigr),\\
[C^{s-2}]\Phi=\beta^{s-2}\bigl(B_d(h)+\beta^2\tau^2\bigr).
\end{gather*}
\end{theorem}
\begin{proof}
$q-2=2s$ and Frobenius give $\Psi=\Phi^2$ pointwise. The binary expansion of $s$ consists of $m-1$ ones, so
$\binom sk\equiv1\pmod2$ for $0\le k\le s$ (Lucas) and $[C^k]\Phi=\sum_az(a)^{s-k}d(a)^k$. Group $a$ in fibres
$\{a,a\oplus\rho\}$ of $d$: $d$ is constant on a fibre, so $[C^s]\Phi=\sum_ad(a)^s=0$. On a fibre,
$L(a\oplus\rho)=L(a)\oplus\tau$, hence $z(a)\oplus z(a\oplus\rho)=\Delta_h(x)\oplus\beta x^2\tau$ with $x=r(a)$
($\mu$ cancels). With $d=\beta x^2$ and $x^{q-1}=1$: $d^{s-1}=\beta^{s-1}x^{-3}$, $d^{s-2}=\beta^{s-2}x^{-5}$
(here $s\ge3$ for $m\ge3$; for $m=3$, $s-2=1$). Therefore
$[C^{s-1}]\Phi=\beta^{s-1}\sum_Hx^{-3}(\Delta_h+\beta x^2\tau)=\beta^{s-1}(A_d+\beta\tau p_1)$ and
$[C^{s-2}]\Phi=\beta^{s-2}\sum_Hx^{-5}(z(a)^2\oplus z(a\oplus\rho)^2)=\beta^{s-2}\sum_Hx^{-5}(\Delta_h+\beta x^2\tau)^2
=\beta^{s-2}(B_d+\beta^2\tau^2p_1)$, and $p_1=1$ (Lemma~\ref{lem:P}).
\end{proof}

\begin{theorem}[Certificate]\label{thm:cert}
If $\ell_d(h)\ne0$ then for every $\FF$-linear $L$ and every $\mu$, $\Phi\ne0$, $\deg\Phi\le q/2-2$ and
\[
\#\{c\in\Fq:\Psi_{L,\mu}(c)\ne0\}\ \ge\ q/2+2 .
\]
More precisely $\deg\Phi=q/2-2$ if $\tau\ne\eta A_d(h)$, and $\deg\Phi=q/2-3$ (so at least $q/2+3$ non-zero
values) if $\tau=\eta A_d(h)$. By Proposition~\ref{dz:prop:link}, the local zero sum of $S_2$ then fails on at
least $q/2+2$ cosets for every local key.
\end{theorem}
\begin{proof}
$\ell_d(h)^2=A_d^2+B_d$. If both coefficients of Theorem~\ref{thm:top} vanished, then $A_d=\beta\tau$ and
$B_d=\beta^2\tau^2=A_d^2$, i.e.\ $\ell_d(h)=0$. So $\Phi\ne0$ and, since $[C^s]\Phi=0$, $\deg\Phi\le s-1$; a non-zero
polynomial of degree $\le s-1$ has at most $s-1$ roots, and $\Psi(c)=0\iff\Phi(c)=0$. If $\tau=\eta A_d$ then the
first coefficient vanishes and the second is $\beta^{s-2}(B_d+A_d^2)=\beta^{s-2}\ell_d^2\ne0$.
\end{proof}
The condition $\ell_d(h)\ne0$ is sufficient; its necessity is not asserted. Theorem~\ref{thm:cert} quantifies over
\emph{all} $\FF$-linear $L$, whereas the local keys of Proposition~\ref{dz:prop:link} only produce the carry-mask maps
$L(a)=a\star\kappa$; the two quantifiers are kept apart in the examples below (Examples~\ref{ex:nonnec}
and~\ref{ex:quant}).

\begin{theorem}[Linearity, explicit correction, exact pass fraction]\label{thm:repair}
Let $\chi:\Fq\to\FF$ be $\FF$-linear with $\chi(\rho)=1$ and $g_d(a)=\chi(a)\,r(a)$. Then
\begin{enumerate}
\item $\ell_d:\Fq^{\Fq}\to\Fq$ is $\Fq$-linear and $\ell_d(g_d)=1$; in particular $\ell_d$ is surjective;
\item $\ell_d(h+\lambda g_d)=\ell_d(h)+\lambda$ for all $h$ and $\lambda\in\Fq$; exactly one of the $q$ functions
$h+\lambda g_d$ fails the certificate, and $h^\sharp=h$ ($\ell_d(h)\ne0$), $h^\sharp=h+g_d$ ($\ell_d(h)=0$) always passes;
\item for every $\Fq$-linear subspace $V\subseteq\Fq^{\Fq}$ with $g_d\in V$, exactly a fraction $1-1/q$ of $V$
passes the certificate; this applies to $V=\Fq^{\Fq}$ and to the space $\mathcal Q_2$ of functions of Boolean
degree $\le2$ (in a fixed basis), $\dim_{\Fq}\mathcal Q_2=1+m+\binom m2$.
\end{enumerate}
\end{theorem}
\begin{proof}
(1) $\Delta_h$ is $\Fq$-linear in $h$; $x\mapsto x^{q/2}$ is additive with $(xy^2)^{q/2}=x^{q/2}y$, so
$B_d(h)^{q/2}=\sum_H(x^{-5})^{q/2}\Delta_h(x)$ is $\Fq$-linear in $h$. Since $r(a\oplus\rho)=r(a)$ and
$\chi(a\oplus\rho)=\chi(a)\oplus1$, $\Delta_{g_d}(x)=x$; hence $A_d(g_d)=p_2=1$, $B_d(g_d)=p_3=0$, $\ell_d(g_d)=1$.
(2) is linearity. (3) $\ell_d|_V$ is a surjective linear functional, so its kernel has index $q$ in $V$.
Each coordinate of $g_d$ is $\chi\cdot r_i$, a product of two affine Boolean functions, so $g_d\in\mathcal Q_2$; a basis of
$\mathcal Q_2$ over $\Fq$ is given by the square-free monomials of degree $\le2$ in the coordinates of $a$.
\end{proof}
The fraction in (3) is the fraction passing a sufficient certificate, not the fraction of $h$ without local zero
sums, and not a fraction of ciphers that are secure.

\subsection{Boolean degree two is minimal for this local problem}\label{ss:deg}
\begin{proposition}\label{prop:affine}
If $h$ is $\FF$-affine then $\Psi_{0,\mu}(c)=0$ for all $c,\mu$. Consequently every $h$ with
$\Psi_{L,\mu}\not\equiv0$ for all $(L,\mu)$ (in particular $g_d$ and every $h$ passing the certificate) has Boolean
degree $\ge2$, and $g_d$ has Boolean degree exactly $2$.
\end{proposition}
\begin{proof}
For $L=0$ and fixed $c,\mu$, $a\mapsto y(a)=h(a)\oplus d(a)c\oplus\mu$ is $\FF$-affine. If it is not injective, every
fibre is a coset of its kernel and has even size, so $\sum_ap(y(a))=0$. If it is bijective,
$\sum_ap(y(a))=\sum_{x\in\Fq}x^{q-2}=\sum_{x}x=0$ ($m\ge2$). By Theorem~\ref{thm:cert}, $\ell_d(g_d)=1$ forces
$\Psi_{0,\mu}\not\equiv0$, so $g_d$ is not affine.
\end{proof}
This is a statement about the Boolean degree of $h$ for this local problem only; it says nothing about AND
counts, about the degree of $F$ or of the whole arrangement, or about optimality among known constructions.

\subsection{Evaluation and the unchanged differential guarantees}
\begin{remark}[Cost of the correction]\label{rem:cost}
The coordinates of $r$ are public affine forms in the bits of $a$ and $\chi(a)$ is one public linear form, so $g_d$
is evaluated with at most $m$ AND gates plus linear gates (in fact $m-1$, Proposition~\ref{prop:mc}); $\lambda g_d$ for a
fixed $\lambda$ adds only an $\FF$-linear map (no AND). This bounds the cost of the correction, not of $h$ itself. The quantities
$\eta,\rho,\chi$ are fixed off line; computing $\ell_d(h)$ for an arbitrary table $h$ may require $q$ evaluations
and is not claimed to be cheap for large $m$. Choosing $h=g_d$ from the start needs no scan of $h$.
\end{remark}
% Finite check: c7/gd (m=5, rho of weight 1 and 3: AND = 5, all 32 points; lambda*g_d as a 5x5 F2 matrix).

\begin{proposition}[Multiplicative complexity of the correction alone]\label{prop:mc}
In the model of XOR/NOT/AND straight-line circuits with two-input AND gates counted and linear gates free,
the multiplicative complexity of $g_d$ (as a map $\FF^m\to\FF^m$) is exactly $m-1$.
\end{proposition}
\begin{proof}
Choose an $\FF$-linear isomorphism $O$ of $\Fq$ whose first coordinate is $\Tr$. Since $r(\Fq)=H$, $Or(a)=(1,v_1(a),\dots,v_{m-1}(a))$
with $v_i$ affine. The map $a\mapsto(t,v_1,\dots,v_{m-1})$, $t=\chi(a)$, is an affine bijection: if two points have the same $v$, their
difference lies in $\{0,\rho\}$ (the fibres of $r$), and $\chi(\rho)=1$ separates them. In these coordinates
$Og_d=(t,tv_1,\dots,tv_{m-1})$, which uses $m-1$ AND gates. Conversely, in a circuit with $k$ AND gates every output is an affine
combination of the inputs and of the $k$ AND outputs, so modulo affine functions the outputs span a space of dimension at most $k$.
The functions $tv_1,\dots,tv_{m-1}$ are linearly independent modulo affine functions (distinct square-free quadratic monomials in
the coordinates $(t,v)$), and an affine change of the input and a linear change of the output do not change this dimension; hence
$k\ge m-1$.
\end{proof}
This concerns $g_d$ computed on its own. It is not a bound for $F$, nor for $h+g_d$, nor for the additional cost when the
correction shares intermediate values with the inverse map (for $h=a^3$ no additional AND gate is needed). The upper bound and
the quotient-dimension lower bound are standard principles of multiplicative complexity; only the identification of the
normal form is specific to $g_d$.
% SOURCE: GPT NEW_DERIVATIONS_JA.md C (proposal), re-proved; c8/checks/p805_p806.py: m=5, d0 (rho=1) and rho=7 (weight 3):
% quadratic rank modulo affine = 4 = m-1, the (m-1)-AND circuit agrees at all 32 points.

\begin{remark}[Same differential guarantees]\label{rem:same}
The correction changes $h$ only. For odd $m\ge5$ the inverse map is an APN permutation, so $(p,d,h^\sharp)$
satisfies Assumption~\ref{ass:PD} and Theorems~\ref{thm:eps}, \ref{thm:W2} (constant $4/N$ when $\wt(\rho)=1$,
$(6+6/q)/N$ in general), \ref{thm:UE} and~\ref{thm:P4} apply to it verbatim, with the same independence and
reset assumptions. For even $m$ only the algebra of Subsections~\ref{ss:norm}--\ref{ss:deg} holds.
\end{remark}

\begin{corollary}[The case $d=d_0$]\label{cor:d0}
For odd $m\ge3$ and $d=d_0$: $\eta=\beta=1$, $\rho=1$, and Theorems~\ref{thm:top}--\ref{thm:repair} reduce to
$[C^{s-1}]\Phi=A+\tau$, $[C^{s-2}]\Phi=B+\tau^2$, $\ell=A+\sqrt B$. For $h(a)=a^3$: $\Delta_h=r$, $A=1$, $B=0$,
$\ell=1$, so $\Psi_{L,\mu}\not\equiv0$ for every local key, with at least $q/2+2$ non-zero values if $\tau\ne1$ and
at least $q/2+3$ if $\tau=1$. The correction $h+\lambda a^3$ and the correction $h+\lambda g_{d_0}$ have the same
effect on $\ell$ ($\ell(a^3)=\ell(g_{d_0})=1$).
\end{corollary}
The function $h=a^3$ remains the main comparison point: it is shared with the inverse chain and its
certificate value is $1$; the general construction does not replace it.

\begin{example}[Vanishing; non-necessity]\label{dz:prop:X21}\label{ex:nonnec}
For $m=5$, $d=d_0$, $h(a)=a^{21}$ and $L=0$, $\mu=0$: $\Psi_{0,0}\equiv0$ on $\mathbb F_{32}$ (substitute $a=b^3$;
the expansion of $(X+C(X^{12}+X^6+1))^{30}$ has $128$ monomials and no positive $X$-degree divisible by $31$).
Here $\ell(h)=0$. Conversely, for random $h$ with $\ell_d(h)=0$ we observed no \emph{carry-mask} key
$(L,\mu)=(a\mapsto a\star\kappa,\mu)$ with $\Psi\equiv0$ (Section~\ref{sec:verif}); so, for the carry-mask key family of this fixed
model, $\ell_d=0$ is not a vanishing criterion. Neither observation is a classification.
\end{example}

\begin{example}[The two quantifiers differ]\label{ex:quant}
Let $m=3$, $\mathbb F_8=\FF[t]/(t^3+t+1)$ with the polynomial basis, $d=d_0$, and let $h$ have the value table
$(h(0),\dots,h(7))=(7,5,2,2,2,0,7,2)$ (index $a=a_0+2a_1+4a_2$). Then $\ell(h)=0$. For each of the $64$ carry-mask keys
$(\kappa,\mu)$ the function $\Psi_{L,\mu}$ with $L(a)=a\star\kappa$ is not identically zero; but for the $\FF$-linear map $L$ with
$(L(1),L(t),L(t^2))=(3,0,0)$ and $\mu=0$ one has $\Psi_{L,0}\equiv0$ (exact computation; among all $512\cdot8$ pairs $(L,\mu)$ with $L$
linear, $64$ give $\Psi\equiv0$). So a finite check over carry-mask keys says nothing about non-necessity for general $L$, and the
statement after Theorem~\ref{thm:cert} is restricted accordingly. This example is not a counterexample to
Theorem~\ref{thm:cert}, which is a sufficient condition.
\end{example}
% SOURCE: GPT NEW_DERIVATIONS_JA.md D (proposal); independent check c8/checks/p805_p806.py (exit 0, c8/results/p805_p806.json).
% SOURCE: c6/h/notes/H_CERT_P6.md:45-48 (X21); c7/gd/notes/P7-05_gd.md (0 -> 0 keys for random ell=0 h).

% sec_5_composition_v4.tex -- C9 paper lane (P9-09/P9-10): internal labels removed from the text ("; C2" after the
% MDS citation, the table name h_{C3}); mathematics unchanged from sec_5_composition_v3.tex.
% sec_5_composition_v3.tex -- C8 (P8-12): evidence kind attached to every attack/cost statement; C7 repair experiment.
% sec_5_composition_v2.tex -- v2: Table tab:fp header; last paragraph corrected (v1 misread P6-09). Group III.
\section{Composition: an averaged bound and what it does not prevent}\label{sec:comp}

This section concerns the arrangement of Section~\ref{ss:arr}, with every local layer built from one
table $F=F_{p,d,h}$ satisfying Assumption~\ref{ass:PD}, $m\ge5$ odd. Everything here is an
application of known principles; we state the results in the form needed and give the short proofs to
keep the paper self-contained, without claiming novelty for the principles.

\subsection{Two standard lemmas}
\begin{lemma}[Post-composition with a reset key; standard]\label{lem:reset}
Let $G$ be a finite abelian group, $E_\kappa:G\to G$ a family of maps, $M$ an automorphism of $G$, and
$C_k:G\to G$ a family of permutations. Let $\kappa$, $L$, $k$ be independent with $L$ uniform on $G$. Then
\[
\E_{\kappa,L,k}\DP_{C_k\circ\tau_L\circ M\circ E_\kappa}(\alpha,\gamma)
=\sum_{\delta}\E_\kappa\DP_{E_\kappa}(\alpha,\delta)\;\bar T_C(M\delta,\gamma),
\]
with $\bar T_C=\E_k T_{C_k}$,
where $\tau_L(y)=y+L$. In particular, if $\E_\kappa\DP_{E_\kappa}(\alpha,\delta)\le b_\star$ for all $\delta$, then the
left side is at most $b_\star$.
\end{lemma}
% SOURCE: c3/notes/MAINLINE_P3.md:24-27 (section D: DDT doubly stochastic, reset count, suffix bound);
%         c4/notes/CORE_P4.md:57 (sum_reset_ddt); Lean: C3 SuffixReset, C4 ResetCompose (Section 7).
% Known principle (same form as Park et al. ASIACRYPT 2002 Lemma 2, c3/notes/MAINLINE_P3.md:27).
\begin{proof}
For a fixed pair $(y,y')$ with $y'-y=M\delta$, the number of $L$ with $C(y'+L)-C(y+L)=\gamma$ is
$\delta_C(M\delta,\gamma)$ (translation). Averaging over the inputs and $\kappa$ gives the identity. Each
$C_k$ is a permutation, so $\bar T_C$ has column sums $1$, and since $M$ is bijective,
$\sum_\delta\bar T_C(M\delta,\gamma)=1$; the bound follows. (Without the reset $L$ the product formula fails in
general.)
\end{proof}

\begin{lemma}[Two-layer MDS count for averaged kernels; standard]\label{lem:D1}
Let $R$ be a finite commutative ring, $M\in\mathrm{Mat}_w(R)$ with all minors units, and let
$T_i,U_j$ ($1\le i,j\le w$) be non-negative $R\times R$ matrices with row sums $1$, $T(0,\cdot)=e_0$,
$T(\alpha,0)=0$ for $\alpha\ne0$, and all entries of non-zero rows at most $b_\star$. Put
$\mathsf K(\alpha,\gamma)=\sum_{\delta\in R^w}\prod_iT_i(\alpha_i,\delta_i)\prod_jU_j((M\delta)_j,\gamma_j)$. Then
$\mathsf K(\alpha,\gamma)\le b_\star^{\,w}$ for every $\alpha\ne0$ and every $\gamma$.
\end{lemma}
% SOURCE: c2/notes/AVERAGED_KERNEL_MDS.md:11-18 (lemma D1); attribution :4 (Hong et al. FSE 2000,
% Park et al., wide-trail/MDS two-layer argument). Not claimed as new.
\begin{proof}
Let $I=\mathrm{supp}\,\alpha$, $J=\mathrm{supp}\,\gamma$, $a=|I|\ge1$, $b=|J|$. Non-zero terms have
$\mathrm{supp}\,\delta=I$ and $\mathrm{supp}\,M\delta=J$. If $a+b\le w$ there is no such $\delta\ne0$ (branch number $w+1$),
so $\mathsf K=0$. If $a+b>w$, choose $I'\subseteq I$ with $|I'|=a+b-w$; the condition $(M\delta)_{J^c}=0$ determines
$\delta_{I\setminus I'}$ from $\delta_{I'}$ because the $(w-b)\times(w-b)$ minor $M_{J^c,I\setminus I'}$ is invertible. Bound the
$w-b$ factors $T_i$, $i\in I\setminus I'$, and the $b$ factors $U_j$, $j\in J$, by $b_\star$, and sum the remaining
$\prod_{i\in I'}T_i(\alpha_i,\delta_i)$ over $\delta_{I'}$ to at most $1$.
\end{proof}

\subsection{The averaged multi-word bound}
\begin{theorem}\label{thm:P4}
Let $q\ge32$, $w\ge2$, $M_1$ with all minors units and $M_2,M_3$ invertible, all key words
independent and uniform. For every $h$ and every $\alpha\ne0$, $\gamma$:
\begin{enumerate}
\item if $\wt(\rho)=1$, then $\E_{\cK}\DP_E(\alpha,\gamma)\le(4/N)^w$;
\item for general $\rho$, $\E_{\cK}\DP_E(\alpha,\gamma)\le((6+6/q)/N)^w\le(99/(16N))^w$.
\end{enumerate}
\end{theorem}
% SOURCE: c5/notes/CORE_P5.md:9-13 (corrected theorem P4, (a)/(b)); proof c4/notes/CORE_P4.md:56-61.
\begin{proof}
Let $b_\star$ be the bound of Theorem~\ref{thm:W2} in the relevant case. By Proposition~\ref{prop:markov}
(with the whitening $k_0$, resp.\ the reset $L_1$, providing the uniform input), the averaged kernel of
each word of $S_2^{(1)}$ and of $S_2^{(2)}$ is $T_B^2T_F$, which satisfies the hypotheses of
Lemma~\ref{lem:D1} with $b_\star$ (it is the averaged kernel of permutations). With $L_1$ uniform and keys
independent across words, the averaged kernel of $x\mapsto X_2$ is the kernel $\mathsf K$ of Lemma~\ref{lem:D1}
for $M=M_1$ (\cite{Hong,Park}), hence at most $b_\star^{\,w}$. Apply Lemma~\ref{lem:reset} with
$(M_2,L_2,S_1^{(3)})$ and then with $(M_3,L_3,S_1^{(4)})$; each $S_1(\cdot;k)$ is a permutation.
\end{proof}

\begin{remark}
The bound is an average over the $10w$ independent key words. It is not a fixed-key bound, it says
nothing about the key schedule, and the reset keys $L_2,L_3$ are essential to the proof. $M_1$ uses the
MDS property; $M_2,M_3$ only invertibility.
\end{remark}

\subsection{A candidate-existence distinguisher}
\begin{lemma}[Zero sum of all ring elements; known]\label{lem:sumR}
$\sum_{r\in R}r=0$ for $R=\GR(4,m)$, $m\ge2$, and $N\cdot y=0$ for every $y\in R$.
\end{lemma}
% SOURCE: c4/notes/CORE_P4.md:30 (premise Z; Knudsen-Wagner FSE 2002 Thm 1 for the general statement).
\begin{proof}
$(R,+)\cong(\mathbb Z/4)^m$; in each coordinate each residue occurs $4^{m-1}$ times and
$4^{m-1}(0+1+2+3)\equiv0\pmod4$ for $m\ge2$. Also $4\mid N$.
\end{proof}

\begin{lemma}[Zero sum at $V_2$]\label{lem:V2}
Fix all keys and let $x$ range over a one-word slice $\{x: x_{a}\in R,\ x_j=c_j\ (j\ne a)\}$. Then
$\sum_xV_2(x)=0$ in $R^w$.
\end{lemma}
% SOURCE: c4/notes/CORE_P4.md:32-33 (proposition Z); c5/notes/CORE_P5.md:48 (lemma, property P).
\begin{proof}
$X_{1,a}$ runs over $R$ bijectively and the other words of $X_1$ are constant. Since $M_{1,ia}$ is a unit,
$V_{1,i}$ is a bijective function of $x_a$ for every $i$, hence so is $X_{2,i}$. By Lemma~\ref{lem:sumR},
$\sum_xX_{2,i}=0$, and $\sum_xV_{2,j}=\sum_iM_{2,ji}\cdot0+N L_{2,j}=0$.
\end{proof}

For a candidate $\theta=(K_4,c,K_3)\in R^w\times R^w\times R^w$ define the word-wise map
\[
G_\theta(y)=S_1^{-1}\bigl(M_3^{-1}S_1^{-1}(y;K_4)-c;\ K_3\bigr).
\]
$G_\theta$ is a bijection of $R^w$, and for $\theta^*=(K_4,M_3^{-1}L_3,K_3)$ one has $G_{\theta^*}(E(x))=V_2(x)$.
The distinguisher $D_s$ queries $s$ pairwise disjoint one-word slices $X_1,\dots,X_s$ (fixed in advance,
non-adaptively; $Q=sN$ distinct inputs) and accepts iff
\[
A:\quad \exists\theta\ \forall t\le s\ \forall j:\ \textstyle\sum_{x\in X_t}G_\theta(y(x))_j=0 .
\]

\begin{theorem}[Candidate-existence distinguisher]\label{thm:V}
Let $M=N^w$ and $Q=sN\le M$.
\begin{enumerate}
\item On the arrangement $E$, $D_s$ accepts with probability $1$ for every key.
\item For a uniformly random function $\RF$, $\Pr[A]\le N^{(3-s)w}$.
\item For a uniformly random permutation $\RP$, $\Pr[A]\le\min\{1,\ N^{(3-s)w}+\binom{Q}{2}/M\}$ and
\[
\Pr\nolimits_{\RP}[A]\le\min\Bigl\{1,\ N^{(3-s)w}\,\frac{M^Q}{(M)_Q}\Bigr\},\qquad (M)_Q=M(M-1)\cdots(M-Q+1).
\]
\end{enumerate}
No independence between different candidates, or between output words, is used.
\end{theorem}
% SOURCE: c5/notes/CORE_P5.md:50-53 (theorem V (1)-(3), additive form);
%         c6/attack/notes/ATTACK_P6_ATK_v2.md:9-28 and c6/notes/resp/P6-07.md:5-13 (conditioned form).
% Attribution: the multiplicative RP/RF comparison is Bernstein 2005 Thm 2.1/2.2 (c6/lit/notes/LIT_P6.md:78-85);
% the additive form is the PRP/PRF switching argument. Neither is claimed as new.
\begin{proof}
(1) Lemma~\ref{lem:V2} and $\theta^*$.
(2) Fix $\theta$. The answers of $\RF$ on the $Q$ distinct inputs are independent and uniform in $R^w$, and
$G_\theta$ is a bijection, so the values $G_\theta(y(x))$ are independent and uniform. The sum over one slice
of $N$ independent uniform elements is uniform, so it vanishes with probability $N^{-w}$; disjoint slices
give independent sums, hence probability $N^{-sw}$. A union bound over the $N^{3w}$ candidates gives the claim.
(3) On fixed distinct inputs, the output distribution of $\RP$ equals that of $\RF$ conditioned on the event
$\mathsf D$ that all $Q$ outputs are distinct, and $\Pr_{\RF}[\mathsf D]=(M)_Q/M^Q$. Hence
$\Pr_{\RP}[A]=\Pr_{\RF}[A\cap\mathsf D]/\Pr_{\RF}[\mathsf D]\le\Pr_{\RF}[A]\,M^Q/(M)_Q$; no independence of $A$ and
$\mathsf D$ is used. The additive form follows from $|\Pr_{\RP}[A]-\Pr_{\RF}[A]|\le\Pr_{\RF}[\overline{\mathsf D}]\le\binom Q2/M$.
\end{proof}

\begin{proposition}[Decomposition of the acceptance condition]\label{prop:dec}
$G_\theta(y)_j$ depends only on $K_4$, $c_j$ and $K_{3,j}$. Hence $A$ holds iff
$\exists K_4\ \bigwedge_j\ \exists(c_j,K_{3,j}):\ P_j(K_4,c_j,K_{3,j})$, where $P_j$ is the slice condition for
word $j$. An implementation that, for each $K_4$, searches each target word separately computes $D_s$
exactly; it examines $N^{w}\cdot N^{2}$ cells per target word.
\end{proposition}
% SOURCE: c5/notes/CORE_P5.md:57 (decomposition, calibration m2 w2 all 16,777,216 theta, m3 sample);
%         c6/attack/notes/ATTACK_P6_ATK_v2.md:44 (N^{w+2} = cells enumerated per target word; not a lower bound).
\begin{proof}
For fixed $K_4$, the conditions for different $j$ involve disjoint variables $(c_j,K_{3,j})$.
\end{proof}

\subsection{Numbers, resources and recovered objects}
Table~\ref{tab:fp} lists the false-acceptance bounds for $\RP$ (exact rational arithmetic).
For two words the additive bound is void, while the conditioned bound is not.
\begin{table}[ht]\small
\caption{Bounds on $\Pr_{\RP}[A]$ from Theorem~\ref{thm:V}(3).}\label{tab:fp}
\begin{tabular}{lrrll}
\toprule
case & $Q$ & $\log_2$ cond. & cond.\ (outward) & additive\\
\midrule
$m=3,w=2,s=4$ & 256 & $-0.2567$ & $<9/10$ & $\approx7.97$ (void)\\
$m=3,w=2,s=5$ & 320 & $-5.536$ & $<1/40$ & $\approx12.46$ (void)\\
$m=5,w=2,s=4$ & 4096 & $-8.446$ & $<1/300$ & $\approx7.998$ (void)\\
$m=5,w=2,s=5$ & 5120 & $-21.94$ & $<1/4\,000\,000$ & $\approx12.50$ (void)\\
$m=5,w=4,s=4$ & 4096 & $-39.99999$ & $<2/N^4$ & $8386561/2^{40}$\\
\bottomrule
\end{tabular}
\end{table}
% SOURCE: c6/attack/notes/ATTACK_P6_ATK_v2.md:30-38; c6/notes/resp/P6-07.md:22-31; m5w4 additive: c5/notes/CORE_P5.md:55

Table~\ref{tab:tdm} fixes, for each experiment, the game, what is unknown, what is recovered, and the
resources; each entry is marked as a theorem, a finite measurement, or an extrapolation. Data is $D=sN$ chosen plaintexts; $T$ is counted in local multiply-accumulate operations of
the implemented enumeration; this is the cost of one implementation, not a lower bound.
\begin{table}[ht]\small
\caption{Resources and recovered objects (the recovered object is at most $\theta$, i.e.\ $3w$ of the $10w$ key words).}\label{tab:tdm}
\begin{tabular}{p{0.15\textwidth}p{0.11\textwidth}p{0.13\textwidth}p{0.21\textwidth}p{0.2\textwidth}}
\toprule
experiment & unknown & recovered & $T/D/M$ & outcome / status\\
\midrule
full-blind, $m=3,w=2$ & all of $\theta$ & $\theta$ & $0.77$\,s, $1.54\cdot10^{10}$ ops; $D=4N=256$; $21$\,MB & $s\ge4$: unique candidate $=$ true $\theta$ in $144/144$ runs; $s=3$ ambiguous in $33/48$ (verified, finite)\\
full-blind, $m=5,w=2$ & all of $\theta$ & --- & extrapolated $2.7\cdot10^{16}$ ops (72--87\,h single process); partial-sum variant $5.6\cdot10^{15}$ ops, $\approx30$\,h (Poisson-type extrapolation from $m=5$ timings: neither a measurement nor a bound) & not run\\
candidate existence, $m=5,w=4,s=4$ & existence of $\theta$ & none (dis\-tin\-guish\-ing only) & $D=4N$; implemented $T\approx5.1\cdot10^{22}$ ops & advantage $\ge1-2^{-39.99}$ (Theorem~\ref{thm:V})\\
\bottomrule
\end{tabular}
\end{table}
% SOURCE: c5/notes/CORE_P5.md:25-33 (contract table); c6/attack/notes/ATTACK_P6_ATK_v2.md:45-53 (T/D/M by s),
%         :74-83 (partial-sum decomposition; the 30 h figure uses p = 1-exp(-N^{2-s}), a Poisson-type
%         approximation, per REVIEW_JA C6-02 -- an extrapolation, neither a measurement nor a bound).
% Controls (m3 w2, PD h0, v2 slices): true 16/16 accepted and unique; RF 0/200; RP 0/200
%         (ATTACK_P6_ATK_v2.md:66-70). These are samples, and s=4/s=5 share output streams.

\subsection{A small averaged differential bound is not PRP security}
The same arrangement satisfies Theorem~\ref{thm:P4} and Theorem~\ref{thm:V}. For $m=5$, $w=4$ the averaged
differential probability is at most $(4/N)^4=2^{-32}$ for every non-zero input difference, while $D_4$
uses $4N=2^{12}$ chosen plaintexts and distinguishes $E$ from a random permutation with advantage at
least $1-2^{-39.99}$. There is no contradiction: Theorem~\ref{thm:P4} is an average over independent keys
of a differential quantity, whereas $D_s$ exploits an integral property (Lemma~\ref{lem:V2}) and a
candidate set of size $N^{3w}$, far smaller than the key space $N^{10w}$. The implemented enumeration
for $m=5$, $w=4$ costs more than the codebook $N^w$; we prove no lower bound on the time of any
attack, and we do not claim that no attack below the codebook exists.

The local correction of Section~\ref{sec:repair} acts inside one $S_2$ block. It removes the local zero
sum of Proposition~\ref{dz:prop:link} for every local key, but it does not remove the word-level zero sum of
Lemma~\ref{lem:V2}, which holds for every bijective local layer. Consistently, in the finite experiments the
tables compared (natural, $h=0$, one further fixed table, $h=a^3$) behaved alike: in the two-layer oracle test all four kept
only the true local candidates ($16/16$ targets), the $m=3$, $w=2$ full-blind runs with $h\in\{0,a^3\}$ were
unique for $s\ge4$, and the SAT experiments on the same arrangement did not separate the tables. With the correction itself
($m=3$, $w=2$, $s=4$, the same $16$ keys and slices for eight tables including $h=a^3$, $h_0$, $g_{d_0}$ and $h_0+g_{d_0}$, the latter
changing $\ell$ from $0$ to $1$ and the number of local keys with $\Psi\equiv0$ from $32$ to $0$), the real construction was accepted in
$16/16$ runs for every table, $\theta$ was recovered uniquely in $127/128$ runs, and $\RF$, $\RP$ were accepted in $0/16$ runs each
(finite measurement; the case $m=5$ was not run). A local certificate and a word-level candidate-existence distinguisher are
different statements; the success of the correction does not neutralize the distinguisher.
% SOURCE: c5/notes/RESPONSE_C5.md:47-49 (P5-07..P5-11); c6/notes/RESPONSE_C6.md:25,28 (P6-09, P6-12).

% sec_6_mixing_v5.tex -- C9 paper lane (P9-01, finding C8-F02): the exponent 1/2 is stated for p(x)=x^{q-2}, odd m,
% d=1+K0(x/rho) in the section opening, the subsection title, the corollary titles and the text before Prop. condsteps;
% internal references ("C6 reduction", "C7 proof", "previous version") removed. Statements and proofs unchanged.
% sec_6_mixing_v4.tex -- C8: exact Z block (P8-02), G block uniform in h and exponent 1/2 for the inverse map (P8-07),
% two-word reduction now proved in Appendix A (P8-04), unconditional step count with the proved exponent (P8-03).
% sec_6_mixing_v3.tex (v3: one display split, layout only) -- written by the main author (C7, P7-06..P7-08, P7-14).
% Proof sources: c7/r1/R1_P7.md (Theorem thm:R1, main author; blocks G and E follow the fourier lane's
% sketches c7/fourier/notes/P7-07_fourier.md C1/C2, re-proved; block Z is new), GPT NEW_DERIVATIONS_JA.md B
% (partial Fourier identity, re-proved) and C (conditional constants, re-proved: c7/a1/notes/P7-08_a1.md),
% C6 c6/a1/notes/N1_P6.md (Proposition D2), C5 c5/a1/notes/LYAP_P5.md (L1-L5), UE_P5.md (1779/256).
\section{Mixing of the alternating arrangement: a uniform one-word contraction}\label{sec:mixing}

This section concerns the averaged (Markov) transition of the alternating two-word arrangement of
\cite{FukuiFFA}-type blocks with an MDS mixing $M$ whose entries and determinant are units. Its statements
are about key-averaged transitions only. Theorem~\ref{thm:R1} gives a contraction exponent that is positive and
uniform over all PD data and all $h$; with it, Theorem~\ref{thm:steps} reaches the scale $1/N$ in a fixed number of steps
without further hypotheses. For the inverse map $p(x)=x^{q-2}$ with $m$ odd and $d=1\oplus K_0(x/\rho)$, $K_0(X)=X^4+X^2$, $\rho\ne0$, the
exponent is $1/2$ (Corollary~\ref{cor:half}), the largest possible for any PD data (Proposition~\ref{prop:Zexact}); for the inverse map with a
general admissible $d$ the constant obtained grows with the linearized degree of $d$, and the exponent $1/2$ uniformly over all PD data is
open (Remark~\ref{rem:exp}). The constants are explicit but weak for small $m$.

\subsection{One word: the partial Fourier picture}
Let $T=T_B$ be the averaged transition matrix of $B$ on ring differences ($N\times N$, doubly stochastic,
$T(\alpha,0)=[\alpha=0]$). For a low digit $u$ of the input difference and $s$ of the output difference, use on
each fibre the orthonormal Walsh basis in the high digit:
\[
\widehat T((u,\eta),(s,\xi))=q^{-1}\sum_{v,t}(-1)^{\langle\eta,v\rangle+\langle\xi,t\rangle}\,T((u,v),(s,t)),
\]
where $\langle\cdot,\cdot\rangle$ is the coordinate inner product of the fixed basis. For $u\notin\{0,\rho\}$ put
$k=K(u)$, $A_a=d(a)d(a\oplus u)/k$, $D_up(a)=p(a\oplus u)\oplus p(a)$ and
$C_a(s)=h(a)\oplus\frac{d(a)}{k}\bigl[s\oplus h(a\oplus u)\oplus h(a)\oplus d(a\oplus u)(a\star u)\bigr]$.

\begin{proposition}[Partial Fourier formulas]\label{prop:pf}
For every $(p,d,h)$ under Assumption~\ref{ass:PD}:
\begin{enumerate}
\item ($\cG$) $\widehat T((u,\eta),(s,\xi))=q^{-2}\sum_{a:\,M_{A_a}^\ast(s\star\xi)=\eta}
(-1)^{\langle\xi,D_up(a)\rangle+\langle s\star\xi,\,C_a(s)\rangle}$, where $M_A:x\mapsto Ax$ and $M_A^\ast$ is its adjoint for
$\langle\cdot,\cdot\rangle$. For $\eta\ne0$ the sum is either empty or runs over one coset $a\oplus\langle\rho,u\rangle$
(four points), and the four signs multiply to $+1$; hence every such entry lies in $\{0,\pm4/q^2\}$.
\item ($\cE$, $u=\rho$) $\widehat T=q^{-2}[s\star\xi=0]\sum_a(-1)^{\langle\eta,v_a(s)\rangle+\langle\xi,D_\rho p(a)\rangle}$ with
$v_a(s)=(s\oplus h(a)\oplus h(a\oplus\rho))/d(a)\oplus(a\star\rho)$.
\item ($\cZ$, $u=0$) $\widehat T=q^{-2}[s\star\xi=0]\sum_a(-1)^{\langle\eta,s/d(a)\rangle}$.
\end{enumerate}
\end{proposition}
\begin{proof}
By Lemma~\ref{lem:diffB}, for each $a$ the high input digit $b$ is determined in the $\cG$ case by the low
equation, and the high output difference is $t=D_up(a)\oplus s\star(A_av\oplus C_a(s))$; summing the characters
over $t$ and then over $v$ leaves the indicator of $\eta=M_{A_a}^\ast(s\star\xi)$. For fixed $z=s\star\xi\ne0$ the map
$A\mapsto M_A^\ast z$ is injective ($\langle z,Ax\rangle=\Tr(\omega Ax)$ for some $\omega\ne0$), so $A_a$ is determined;
$A_a$ is constant on $a\oplus\langle\rho,u\rangle$ and the two roots $\{y,y\oplus k\}$ of $y^2+ky+kA=0$ have
$d$-preimages forming one such coset. On that coset $D_up$ takes one value twice in each of two pairs, and the
four values of $C_a(s)$ sum to $0$ (the $h$-terms cancel because $d(a)/k+d(a\oplus u)/k=1$, the $s$-terms and the
$\star$-terms appear with equal coefficients an even number of times). In the $\cE$ and $\cZ$ cases the low equation
does not involve $b$; the sum over $b$ of $(-1)^{\langle s\star\xi,\,h(a)\oplus d(a)b\rangle}$ equals $q[s\star\xi=0]$, and
the low equation determines $v$ from $(a,s)$.
\end{proof}
The entrywise bound $4/q^2$ is not an operator-norm bound; the next theorem bounds the operator norm.
% Finite check: author DDT -> integer Walsh, all components, m=3 (24,576) and m=5 (7,864,320), 16 tables
% including rho of weight 1,2,3 (c7/fourier/notes/P7-06_fourier.md); also GPT checks/partial_fourier.py.

\subsection{The uniform contraction R1}
Let $W=R\setminus\{0\}$, $T'=T|_{W\times W}$, and let $\perp\subset\ell^2(W)$ be the functions whose sum over every
low-digit fibre $f_s\cap W$ vanishes; $\tau:=\|\Pi_\perp T'^{\top}\Pi_\perp\|$ (the one-word fluctuation contraction used in
the two-word reduction of Appendix~\ref{app:mix}).

\begin{theorem}[R1]\label{thm:R1}
For every $m\ge3$, every $(p,d,h)$ under Assumption~\ref{ass:PD}:
\[
\tau^2\ \le\ \frac{4(1.5^m-1)}{q}+\frac{3}{q}+\frac{3^m-2q+1}{q^2}\ =\ 5\Bigl(\frac34\Bigr)^m-\frac3q+\frac1{q^2},
\]
so that $\tau\le\sqrt5\,q^{-\gamma}$ with $\gamma=1-\tfrac12\log_23\approx0.2075$.
\end{theorem}
\begin{proof}
\emph{Blocks.} The input space $\perp$ is the orthogonal sum of the fibres $u\ne0$ (Walsh directions $\eta\ne0$)
and the fibre $u=0$ (zero-sum functions on $q-1$ points, extended by $0$ at the origin: an isometry into the
directions $\eta\ne0$). Hence $\Pi_\perp T'^\top\Pi_\perp=[M_\cG\,|\,M_\cE\,|\,M_\cZ]$ and
$\tau^2\le\|M_\cG\|^2+\|M_\cE\|^2+\|M_\cZ\|^2$. On an output fibre $s\ne0$, $\perp$ is spanned by $\xi\ne0$; on the
output fibre $s=0$ ($q-1$ points), the projection of $f|_W$ onto zero-sum functions has norm at most
$\|f-\mathrm{mean}(f)\|$, the norm of the $\xi\ne0$ part of $f$. So each $\|M_\bullet\|$ is at most the norm of the
corresponding rows $\eta\ne0$ and columns $\xi\ne0$ of $\widehat T$.

\emph{$\cG$.} Let $g=q^2\widehat T$ on these rows and columns and $Q=gg^\top$. For rows $r=(u,\eta)$, $r'=(u',\eta')$,
atoms $a\in$ fibre of $r$ and $a'\in$ fibre of $r'$ share the column $(s,\xi)$ iff $s\star\xi=z$, where $z=z_a(\eta)$ is the
unique solution of $M^\ast_{A_a}z=\eta$, and $z_{a'}(\eta')=z$; then $z\ne0$. Summing over the $2^{m-\wt s}$ vectors $\xi$
with $s\star\xi=z$ gives $2^{m-\wt s}(-1)^{\langle z,\delta\rangle}[\delta\subseteq s]$ with
$\delta=D_up(a)\oplus D_{u'}p(a')$. The function $h$ enters only through the phase $\langle z,C_a(s)\oplus C'_{a'}(s)\rangle$;
taking absolute values and summing over $s\supseteq z\vee\delta$, $\sum_{s\supseteq w}2^{m-\wt s}=3^{m-\wt w}$, gives
$|Q[r,r']|\le\sum_{(a,a')}3^{m-\wt(z\vee\delta)}$. For fixed $r$ and $a$, every pair $(u',a')$ determines $\eta'$, so the
absolute row sum is at most $\sum_a\sum_{u',a'}3^{m-\wt(z_a\vee\delta)}$. As $p$ is a permutation, each value
$y=D_{u'}p(a')$ is taken at most $q$ times, so the inner sum is at most
$q\sum_{\delta\in\Fq}3^{m-\wt(z\vee\delta)}=q\,4^m2^{-\wt z}$. Each $z\ne0$ arises from at most four $a$ (one coset), so
$\sum_a2^{-\wt z_a}\le4((3/2)^m-1)$. By Schur's test $\|M_\cG\|^2\le\|Q\|/q^4\le4(1.5^m-1)/q$. ($\cG$ rows do not reach
the output fibre $s=0$: there $s\star\xi=0$ and $\eta\ne0$ leaves an empty sum.)

\emph{$\cE$.} Write the block as $q^{-2}\Phi N$ with $\Phi[\eta,v]=(-1)^{\langle\eta,v\rangle}$ ($\|\Phi\|\le\sqrt q$) and
$N[v,(s,\xi)]=[s\star\xi=0]\sum_{a:v_a(s)=v}(-1)^{\langle\xi,D_\rho p(a)\rangle}$. In $NN^\top[v,v']$ the sum over $\xi\ne0$
with $s\star\xi=0$ is at most $2^{m-\wt s}[\delta\subseteq s]+1$ in absolute value. Since $D_\rho p$ is two-to-one ($p$ APN),
for fixed $a$ at most $2^{\wt s+1}$ values $a'$ have $\delta\subseteq s$; so the row sum is at most
$\sum_s\#\{a:v_a(s)=v\}\cdot3q=3q^2$, because $s\mapsto v_a(s)$ is bijective for each $a$. Hence
$\|M_\cE\|^2\le q^{-4}\cdot q\cdot3q^2=3/q$.

\emph{$\cZ$.} Here $s=d(a)v\ne0$, and the entry does not depend on $\xi$: for fixed $s$ the block has
$n_s=2^{m-\wt s}-1$ equal columns $x_s[\eta]=q^{-2}\widehat w_s(\eta)$, $w_s=2\cdot\mathbf 1_{s(\Img d)^{-1}}$
($d$ is two-to-one onto a set of size $q/2$). By Parseval, $\sum_\eta\widehat w_s(\eta)^2=2q^2$ and $\widehat w_s(0)=q$,
so $\|x_s\|^2=q^{-2}$ on $\eta\ne0$, and $\|M_\cZ\|^2\le\sum_{s\ne0}n_s\|x_s\|^2=(3^m-2^m-q+1)/q^2$.
\end{proof}
The function $h$ enters only the phases of the $\cG$ and $\cE$ blocks; the $\cZ$ block depends only on $\Img d$.
The $\cG$ estimate uses that $p$ is a permutation and that $d$ is affine; the $\cE$ estimate uses APN in the direction
$\rho$; oddness of $m$ is not used.
% Finite check: m=5, 8 tables: tau 0.139-0.151 (bound vacuous); blocks ||M_G||^2 <= 0.0057 (bound 0.82),
% ||M_E||^2 <= 0.022 (bound 0.094), ||M_Z||^2 = 15/1024 exactly (bound 180/1024); m=7: tau 0.063 (bound 0.80).


\subsection{The exact $\cZ$ block}
\begin{proposition}[Exact $\cZ$ norm]\label{prop:Zexact}
For every $m\ge3$ and all $(p,d,h)$ under Assumption~\ref{ass:PD}, $\|M_\cZ\|^2=(q/2-1)/q^2$. Consequently
\[
\tau^2\ \le\ 4\Bigl(\frac34\Bigr)^m-\frac1{2q}-\frac1{q^2}.
\]
\end{proposition}
\begin{proof}
By Lemma~\ref{lem:norm} (which uses only (P2)), $\Img d=\{z:\Tr(\eta_0z)=1\}$ for a unique $\eta_0\ne0$. For $s\ne0$ let
$w_s(v)=2\cdot[\Tr(\eta_0s/v)=1]$ for $v\ne0$ and $w_s(0)=0$, and $x_s(\nu)=q^{-2}\sum_v(-1)^{\langle\nu,v\rangle}w_s(v)$ for $\nu\ne0$; by
the proof of Theorem~\ref{thm:R1} these are the columns of the $\cZ$ block, each repeated $n_s=2^{m-\wt(s)}-1$ times.
One has $\sum_vw_s(v)=q$, $\sum_vw_s(v)^2=2q$, and for $s\ne t$, after the substitution $v\mapsto1/v$, the conditions
$\Tr(\eta_0sv)=1$ and $\Tr(\eta_0tv)=1$ are two independent affine conditions (the linear forms differ), so
$\sum_vw_s(v)w_t(v)=4\cdot q/4=q$. By Parseval over all $\nu$ and removal of $\nu=0$,
\begin{gather*}
\langle x_s,x_t\rangle=q^{-4}\Bigl(q\sum_vw_sw_t-\bigl(\textstyle\sum_vw_s\bigr)\bigl(\sum_vw_t\bigr)\Bigr)=q^{-2}[s=t],\\
\sum_{\nu\ne0}x_s(\nu)=q^{-2}\bigl(qw_s(0)-\textstyle\sum_vw_s(v)\bigr)=-\frac1q .
\end{gather*}
Hence on the extended space of all functions of $\nu\ne0$, $M_\cZ^*M_\cZ=\sum_{s\ne0}n_sx_sx_s^{\top}$ has the orthogonal
eigenvectors $x_s$ with eigenvalues $n_s/q^2$, the largest being $(q/2-1)/q^2$ ($\wt(s)=1$); this gives
$\|M_\cZ\|^2\le(q/2-1)/q^2$. The physical input space of the $\cZ$ block is the space of functions on the $q-1$ points $(0,v)$,
$v\ne0$, with sum $0$; extended by $0$ at the origin, its image under the high-digit Walsh transform is exactly the set of
coefficient vectors $c$ with $c(0)=0$ and $\sum_{\nu\ne0}c(\nu)=0$ (the value at the origin is $q^{-1/2}\sum_\nu c(\nu)$). For two
distinct $s,t$ of weight one ($m\ge2$), $c=x_s-x_t$ satisfies both conditions and is an eigenvector for $(q/2-1)/q^2$. On these
inputs the output low digit is $s=d(a)v\ne0$, so the output fibre $s=0$ (where the proof of Theorem~\ref{thm:R1} only uses a contraction) is not
reached, and on the fibres $s\ne0$ the projection onto $\perp$ is exactly the removal of $\xi=0$. So the bound is attained by
a physical input, and $\|M_\cZ\|^2=(q/2-1)/q^2$. The bound on $\tau^2$ is Theorem~\ref{thm:R1} with the $\cZ$ term replaced:
$4(1.5^m-1)/q+3/q+(q/2-1)/q^2=4(3/4)^m-1/(2q)-1/q^2$.
\end{proof}
% SOURCE: GPT NEW_DERIVATIONS_JA.md A (audit proposal), re-proved here (c8/notes/resp/P8-02.md);
% finite check c8/checks/p802_zexact.py: m=3,5, d0 and d=1+K0(x/rho) (rho=7,11), h in {0,a^3,random}: integer Walsh of the
% physical Z rows = x_s columns, exact Gram q^2[s=t], SVD norm = (q/2-1)/q^2, exact Rayleigh quotient of the witness.
Since $\tau\ge\|M_\cZ\|=((q/2-1)/q^2)^{1/2}$, no bound $\tau=O(q^{-\gamma'})$ with $\gamma'>1/2$ can hold uniformly; $1/2$ is the
largest possible exponent.

\subsection{The $\cG$ block: removal of $h$, and exponent $1/2$ for the inverse map with $d$ of $d_0$ type}
For $u\notin\{0,\rho\}$, $k=K(u)$ and $a\in\Fq$ put $\mu_{u,a}=d(a)/k$, $y_{u,a}=D_up(a)$, and for $z\ne0$ let
$\ell_{u,a}(z)=M^*_{\mu_{u,a}}z$. Let $g=q^2\widehat T$ on the rows $(u,\eta)$, $u\notin\{0,\rho\}$, $\eta\ne0$, and the columns $(s,\xi)$,
$\xi\ne0$; by Proposition~\ref{prop:pf} and the proof of Theorem~\ref{thm:R1}, $\|M_\cG\|=\|g\|/q^2$ (the $\cG$ rows do not reach $s=0$).
Define
\[
\mathsf M(p,d)=\max_{y,\nu\in\Fq}\#\bigl\{(u,a):u\notin\{0,\rho\},\ D_up(a)=y,\ d(a)=\nu K(u)\bigr\}.
\]

\begin{theorem}[$\cG$ block, uniform in $h$]\label{thm:G}
For every $m\ge3$, all $(p,d)$ under Assumption~\ref{ass:PD} and every $h$,
$\|M_\cG\|^2\le\mathsf M(p,d)/q$.
\end{theorem}
\begin{proof}
\emph{Separation of $h$.} In Proposition~\ref{prop:pf}(1) write $C_a(s)=h(a)\oplus\mu_{u,a}(s\oplus E_{u,a})$ with
$E_{u,a}=h(a\oplus u)\oplus h(a)\oplus d(a\oplus u)(a\star u)$. For a column with $s\star\xi=z$,
$\langle z,C_a(s)\rangle=\langle z,h(a)\oplus\mu_{u,a}E_{u,a}\rangle\oplus\langle\ell_{u,a}(z),s\rangle$, so the sign of the atom $(u,a)$ is
$\varepsilon_{u,a}(z)\,\sigma^0_{u,a,z}(s,\xi)$ with $\varepsilon_{u,a}(z)=\pm1$ independent of $(s,\xi)$ and
$\sigma^0_{u,a,z}(s,\xi)=(-1)^{\langle\xi,y_{u,a}\rangle+\langle s,\ell_{u,a}(z)\rangle}$ independent of $h$.

\emph{Factorization.} Let $\Psi$ be the $0/1$ matrix with rows $(u,\eta)$ and columns $(u',a,z)$, $z\ne0$, entry
$[u=u'][\eta=M^*_{A_{u,a}}z]$, and $\Sigma$ the matrix with rows $(u,a,z)$ and columns $(s,\xi)$, entry
$[s\star\xi=z]\,\varepsilon_{u,a}(z)\sigma^0_{u,a,z}(s,\xi)$. Proposition~\ref{prop:pf}(1) says $g=\Psi\Sigma$ (columns with $s\star\xi=0$ vanish on
both sides because $\eta\ne0$). Since $A_{u,a}=d(a)d(a\oplus u)/k\ne0$, $M^*_{A_{u,a}}$ is invertible, so for each row $(u,\eta)$ and each
$a$ exactly one $z$ occurs, and different rows have disjoint supports: $\Psi\Psi^\top=qI$, $\|\Psi\|^2=q$. Every column $(s,\xi)$ of
$\Sigma$ belongs to the single value $z=s\star\xi$, so $\Sigma\Sigma^\top$ is block diagonal over $z$, with blocks
$\mathcal B_z=E_z\mathcal B^0_zE_z$, $E_z=\mathrm{diag}(\varepsilon_{u,a}(z))$, $\mathcal B^0_z=\Sigma^0_z\Sigma^{0\top}_z$. As $E_z$ is a signed
identity, $\|\mathcal B_z\|=\|\mathcal B^0_z\|$, and $\|g\|^2\le\|\Psi\|^2\|\Sigma\|^2=q\max_{z\ne0}\|\mathcal B^0_z\|$; $h$ has disappeared.

\emph{Frame bound.} The non-zero spectrum of $\mathcal B^0_z$ equals that of the sum over all atoms $(u,a)$ of the rank-one
operators $\sigma^0_{u,a,z}\sigma^{0\top}_{u,a,z}$ on the columns
$\{(s,\xi):s\star\xi=z\}$. Each $\sigma^0_{u,a,z}$ is the restriction of a character of $\FF^{2m}$, namely
$\chi_P(\xi,s)=(-1)^{\langle\xi,y\rangle+\langle s,\ell\rangle}$ with $P=(y,\ell)=(y_{u,a},\ell_{u,a}(z))$. If $n_z(P)$ atoms share the point $P$, then, in the semidefinite
order,
\[
\sum_{(u,a)}\sigma^0\sigma^{0\top}=\sum_Pn_z(P)\chi_P\chi_P^\top\preceq\bigl(\max_Pn_z\bigr)\sum_{P\in\FF^{2m}}\chi_P\chi_P^\top
=\bigl(\max_Pn_z\bigr)\,q^2I
\]
by orthogonality of characters. For $z\ne0$ the map $\nu\mapsto M^*_\nu z$ is $\FF$-linear and injective, so $\ell_{u,a}(z)$ determines $\mu_{u,a}$ and
$\max_Pn_z(P)=\mathsf M(p,d)$. Thus $\|\mathcal B^0_z\|\le\mathsf M q^2$, $\|g\|^2\le\mathsf Mq^3$ and $\|M_\cG\|^2=\|g\|^2/q^4\le\mathsf M/q$.
\end{proof}
The absolute-value Schur estimate in the proof of Theorem~\ref{thm:R1} loses the factor $(3/2)^m$ in the cross terms between
different $u$; the proof above avoids it by passing from the row Gram matrix to the atom Gram matrix, where $h$ acts by a diagonal
sign change. The theorem uses of the assumption only $d\ne0$ and the definition of $K$.

\begin{proposition}[Multiplicity for the inverse map]\label{prop:Minv}
Let $p(x)=x^{q-2}$ and $K(x)=\sum_{j=0}^ek_jx^{2^j}$ with $k_e\ne0$. Then $\mathsf M(p,d)\le2^{e+1}+2$. In particular
$\mathsf M\le10$ for $d=d_0$ and, more generally, for $d(x)=1\oplus K_0(x/\rho)$ with $K_0(x)=x^4+x^2$ and any $\rho\ne0$, $m$ odd.
\end{proposition}
\begin{proof}
Fix $(y,\nu)$. As $p$ is a permutation, $y=D_up(a)$ with $u\ne0$ forces $y\ne0$, and for given $a$ at most one $u$ works. For $a\notin\{0,1/y\}$
the equation $1/(a+u)+1/a=y$ gives $u=a^2y/(1+ay)$. Multiplying $\varepsilon\oplus K(a)=\nu K(u)$ by $(1+ay)^{2^e}\ne0$ gives a polynomial
equation $\Phi(a)=0$ of degree at most $2^{e+1}$ (the term $k_j(a^2y)^{2^j}(1+ay)^{2^e-2^j}$ has degree $2^e+2^j$), and $\Phi(0)=\varepsilon\ne0$
because $0\in\Img K$ and $\varepsilon\notin\Img K$. Adding the two exceptional values of $a$ gives the bound. For $d_0=1\oplus K_0$, $e=2$; the
same holds for $K_0(x/\rho)$, whose kernel is $\{0,\rho\}$ and whose image is $\Img K_0=\{\Tr=0\}\not\ni1$ for odd $m$.
\end{proof}
% SOURCE: c8/r1/P8-07_G.md (Theorem G, Prop. I; main author C8 r1 lane), independent adversarial audit c8/r1/P8-07_AUDIT.md
% (all claims CONFIRMED; own brute-force DDT -> Walsh, m=3: 16 tables, m=5: 21 tables; exact M for p=inverse, m=3..11, ~150 cases,
% max M = 5 for d0 type; adversarial h: ||M_G||^2 q/M <= 0.064 at m=5).

\begin{corollary}[Exponent $1/2$ for the inverse map with $d=1\oplus K_0(x/\rho)$]\label{cor:half}
For $p(x)=x^{q-2}$, $m\ge3$ odd, $d=1\oplus K_0(x/\rho)$ (any $\rho\ne0$; $\rho=1$ is $d_0$) and every $h$:
$\tau^2\le13/q+(q/2-1)/q^2<13.5/q$, that is $\tau<3.68\,q^{-1/2}$. For general PD data,
$\tau^2\le(\mathsf M(p,d)+3)/q+(q/2-1)/q^2$.
\end{corollary}
\begin{proof}
$\tau^2\le\|M_\cG\|^2+\|M_\cE\|^2+\|M_\cZ\|^2$ (proof of Theorem~\ref{thm:R1}), Theorem~\ref{thm:G}, the bound $\|M_\cE\|^2\le3/q$, and
Proposition~\ref{prop:Zexact}.
\end{proof}
For general APN permutations $p$ (or $K$ of large $2$-degree) a uniform bound on $\mathsf M(p,d)$ is not proved; in the computed cases
($p\in\{x^{-1},x^3,x^5,x^{13}\}$, random $K$, $m\le11$) $\mathsf M\le11$. Uniformity of the exponent $1/2$ over all PD data is therefore
reduced to the $h$-free inequality $\max_{z\ne0}\|\mathcal B^0_z\|\le Cq^2$ (Remark~\ref{rem:exp}).

\begin{remark}[What limits the exponent]\label{rem:exp}
For general PD data, Theorem~\ref{thm:R1} gives the exponent $\gamma=1-\frac12\log_23\approx0.2075$, and
Theorem~\ref{thm:G} improves it to $1/2$ whenever $\mathsf M(p,d)$ is bounded. Numerically $\|M_\cG\|\approx2.2/q$ ($m=3,5,7$),
$\|M_\cE\|\approx0.8\,q^{-1/2}$ and $\|M_\cZ\|^2=(q/2-1)/q^2$ exactly. What remains open for exponent $1/2$ uniformly over all PD data is
the large-sieve type bound $\max_{z\ne0}\|\mathcal B^0_z\|\le Cq^2$ with an absolute $C$; a bound $\sup\mathsf M<\infty$ suffices.
\end{remark}

\subsection{Two words}
Let $M=\bigl[\begin{smallmatrix}m_{11}&m_{12}\\m_{21}&m_{22}\end{smallmatrix}\bigr]$ have unit entries and unit determinant. The
\emph{alternating arrangement} A1 applies, in every step, $B$ to each word with an independent uniform key in front, and then $M$.
Its averaged transition on two-word differences is $\mathcal TX=\mathrm{Perm}_M(T^\top XT)$, where $X$ is a function on $R^2$ (an
$N\times N$ array) and $(\mathrm{Perm}_MX)(M\beta)=X(\beta)$ (Proposition~\ref{prop:markov} applied word-wise; Section~\ref{ss:arr}
is a different arrangement). Starting from $\pi_0=e_\alpha$, $\alpha\ne(0,0)$, put $\pi_j=\mathcal T^j\pi_0$ and $D_j=\pi_j-U$, $U$ the uniform
distribution on $R^2\setminus\{(0,0)\}$. The spaces $P$, $Q$, the quantities $\theta_2=\|Q\mathcal T^2Q\|$, $\zeta_2=\|Q\mathcal T^2P\|$,
$\eta_{\mathrm{rss}}=\sqrt{2/q+1/(q^2-1)}$ and the proof of the following reduction are in Appendix~\ref{app:mix}
(Proposition~\ref{prop:appD2}): with $\nu^2\le2/(q-1)+\tau^2$,
\begin{equation}\label{eq:D2}
\theta_2\ \le\ \eta_{\mathrm{rss}}+\sqrt{\bigl(\nu^2+(2+\sqrt2)/\sqrt{q-1}\bigr)^2+2\nu^2}.
\end{equation}
% SOURCE: c6/a1/notes/N1_P6.md sections 1-5, now Appendix A (c8 P8-04).
\begin{corollary}[Unconditional two-word contraction]\label{cor:N1}
For all $m\ge3$, all PD data and all $h$, \eqref{eq:D2} with Theorem~\ref{thm:R1} gives
$\theta_2=O(q^{-\gamma})$; with $x=(\sqrt3/2)^m$, $\theta_2\le(1919/240)\,x<8x$ for $m\ge5$ (Lemma~\ref{lem:app-arith}).
\end{corollary}

\begin{theorem}[Step count with the proved exponent]\label{thm:steps}
Let $m\ge5$, all PD data, all $h$, and A1 as above. With $x=(\sqrt3/2)^m=q^{-\gamma}$, for every integer $k\ge0$,
\[
\|D_{2+2k}\|_2\ \le\ \frac85\,10^k\,q^{-1}x^{k-1}.
\]
In particular $\|D_{14}\|_2\le1\,600\,000\,q^{-(1+5\gamma)}\le1\,600\,000/N$, and for each word $j$,
\[
\Pr\nolimits_{\pi_{15}}[\beta_j=0]\ \le\ \min\Bigl(1,\ \frac{11\,118\,751}{N}\Bigr).
\]
\end{theorem}
\begin{proof}
\emph{Start.} By Theorem~\ref{thm:eps} every entry of a non-zero row of $T_B$ is at most $1/q$. Lemma~\ref{lem:D1} with $w=2$, $T_i=U_j=T_B$,
$b_\star=1/q$, followed by the permutation $\mathrm{Perm}_M$, gives $\|\pi_2\|_\infty\le1/N$, hence $\|D_2\|_2\le\|\pi_2\|_2\le(\|\pi_2\|_\infty\|\pi_2\|_1)^{1/2}=1/q$
($D_2\perp U$). For the cell projection, $\|PD\|_2^2\le\sum_C m_C^2/|C|\le1/(q^2-1)$, where $m_C$ is the mass of $\pi_2$ in the cell $C$ and every
cell has at least $q^2-1$ points. Thus $a_2:=\|PD_2\|\le(q^2-1)^{-1/2}$ and $b_2:=\|QD_2\|\le1/q$.

\emph{Comparison.} By Lemma~\ref{lem:app-cmp}, $(a_{j+2},b_{j+2})^\top\le A_q(a_j,b_j)^\top$ entrywise with
$A_q=\bigl[\begin{smallmatrix}\eta_{\mathrm{rss}}+q^{-4}&\eta_{\mathrm{rss}}(1+q^{-2})\\ \zeta_2&\theta_2\end{smallmatrix}\bigr]$, $\zeta_2\le1$. Put $u=q^{-1/2}$.
By Lemma~\ref{lem:app-arith}, $A_q\le\bigl[\begin{smallmatrix}2u&2u\\1&8x\end{smallmatrix}\bigr]$ and, with $\Delta=\mathrm{diag}(x,1)$, the column sums of
$\Delta^{-1}\bigl[\begin{smallmatrix}2u&2u\\1&8x\end{smallmatrix}\bigr]\Delta$ are $2u+x$ and $2u/x+8x$, both at most $10x$ (since $u\le x^2$). For the weighted quantity
$v_j=a_j/x+b_j$ (the $\ell^1$ norm of $\Delta^{-1}(a_j,b_j)^\top$) this gives $v_{j+2}\le10x\,v_j$.

\emph{Conclusion.} $v_2\le[(1-q^{-2})^{-1/2}+x]/(qx)\le8/(5qx)$ because $(1-q^{-2})^{-1/2}\le11/10$ and $x\le1/2$. As $x\le1$,
$\|D_j\|_2\le a_j+b_j\le v_j$, and $v_{2+2k}\le(10x)^kv_2$ gives the general bound. For $k=6$, $x^5\le1/q$ because $3^5<2^8$; hence
$\|D_{14}\|\le(8/5)10^6q^{-2}$. Finally Lemma~\ref{lem:app-line} gives
$|\Pr_{\pi_{15}}[\beta_j=0]-1/(N+1)|\le(1779/256)\|D_{14}\|_2$, and $1/(N+1)+(1779/256)\cdot1\,600\,000/N\le11\,118\,751/N$.
\end{proof}
% SOURCE: GPT NEW_DERIVATIONS_JA.md B (audit proposal), re-proved (c8/notes/resp/P8-03.md); exact rational checks of every
% elementary inequality for m=5..61 and the m=3 start (all 4095 input pairs) in c8/checks/p803_steps.py (exit 0).
The constants are far from sharp and the theorem is not a recommendation of a number of rounds for any small $m$; it is a statement,
uniform in $h$ and in the PD data, that a fixed number of steps reaches the scale $1/N$. It concerns the key-averaged difference
distribution of A1, not the $(2,2,1,1)$ arrangement of Section~\ref{sec:comp}, and nothing here concerns a fixed key, a key schedule or
PRP security.

For the inverse map with $m$ odd and $d=1\oplus K_0(x/\rho)$, Corollary~\ref{cor:half} supplies the hypothesis of the following proposition.

\begin{proposition}[Step count under $\theta_2\le C q^{-1/2}$]\label{prop:condsteps}
Let $q\ge32$, $u=q^{-1/2}$, and assume $\theta_2\le Cu$. With $A_q$ as above one has
$0\le A_q\le B(u)=\bigl[\begin{smallmatrix}2u&2u\\1&Cu\end{smallmatrix}\bigr]$ entrywise,
$\|B(u)^2\|_1\le K_Cu$ with $K_C=\max(C+24/5,\ 14/5+2C/5+C^2/5)$, hence $\|A_q^6\|_1\le K_C^3q^{-3/2}$,
$\|D_{13}\|_2\le\sqrt2K_C^3/N$, and the probability that one word is zero after the next step is at most
$[1+(1779/256)\sqrt2K_C^3]/N$.
\end{proposition}
\begin{proof}
$\eta_{\mathrm{rss}}\le\frac32u$ is equivalent to $q^2-4q-1\ge0$; $q^{-4}\le u/2$; $\frac32(1+q^{-2})\le2$; and $u\le1/5$ for $q\ge25$.
The column sums of $B^2$ are $u(4u+4+C)\le(C+24/5)u$ and $u(2+(4+2C+C^2)u)\le(14/5+2C/5+C^2/5)u$. For
entrywise non-negative matrices $0\le A\le B$ implies $0\le A^k\le B^k$, and the maximal column sum is monotone and
submultiplicative. By Theorem~\ref{thm:eps} and Lemma~\ref{lem:app-start}, $\|D_1\|_2\le q^{-1/2}$ and $a_1+b_1\le\sqrt2\|D_1\|$; so
$\|D_{1+2k}\|_2\le\sqrt2q^{-1/2}\|A_q^k\|_1$, used with $k=6$, and Lemma~\ref{lem:app-line} converts the $\ell^2$ bound into a zero-word
probability ($m\ge5$).
\end{proof}
% SOURCE: c7/a1/notes/P7-08_a1.md (GPT C6 NEW_DERIVATIONS C, re-proved in C7); c6/a1/notes/LYAP_P6.md section 5.

\begin{corollary}[Inverse map with $d=1\oplus K_0(x/\rho)$: thirteen steps]\label{cor:inv13}
For $p(x)=x^{q-2}$, odd $m\ge5$, $d=1\oplus K_0(x/\rho)$ and every $h$, $\theta_2\le(489/50)\,q^{-1/2}$; hence
$\|D_{13}\|_2\le24\,405/N$ and $\Pr_{\pi_{14}}[\beta_j=0]\le169\,597/N$.
\end{corollary}
\begin{proof}
Insert $\tau^2\le13/q+(q/2-1)/q^2$ (Corollary~\ref{cor:half}) and $\nu^2\le2/(q-1)+\tau^2$ in \eqref{eq:D2}. Divided by $u$, each of the terms
$\eta_{\mathrm{rss}}/u$, $(\nu^2+(2+\sqrt2)/\sqrt{q-1})/u$ and $\nu/u$ is non-increasing in $q$ for $q\ge32$, so the ratio $\theta_2/u$ is at most its value at
$m=5$, which is below $489/50$ (outward rational evaluation). Then $K_C<25.85$ for $C=489/50$, and Proposition~\ref{prop:condsteps} gives
$\sqrt2K_C^3<24\,405$ and $1+(1779/256)\sqrt2K_C^3<169\,597$.
\end{proof}
% SOURCE: c8/checks/p803_steps.py part B (ratio list m=5..201, K_C = 25.84168, D13 coefficient 24404.91, zero-word coefficient 169596.07).

% sec_7_verification_v12.tex = v11 with the wording of Lean verification changed to "checked in Lean" / "checked by the Lean kernel" (5 places)
% and the m = 3, m = 5 entries of the adversarial search in Table tab:finite2 stated as the largest observed floating-point
% values (0.46287, 0.06390) instead of the bound "<= 0.46" (which the recorded maximum 0.46286958505344594 exceeds).
% sec_7_verification_v11.tex = v10 with the middle column of the three Lean tables 0.41 -> 0.40\textwidth
% (anonymous iacrtrans build: three 1.05pt overfull boxes); no change of content.
% v10 = v9 + \clearpage at the end so that Tables 4-9 are placed after Section 7
% (in v9 they all floated behind the references). v9 = v8 with the Lean tables placed [!htbp] (v8: Table 8 floated to the end). v8 = v7 with the last Lean table split in two (v7: float too large by 266pt in amsart).
% v7: formalization subsection rewritten for the additive two-word mixer;
% internal session labels and provenance comments removed.
\section{Finite calibration, formalization and cost: the actual scope}\label{sec:verif}

The theorems of Sections~\ref{sec:diff}--\ref{sec:mixing} are proved on paper. This section records
which parts were additionally checked on finite instances, which parts are checked by the Lean kernel,
and what was measured about cost. A finite check is a calibration, not a proof of a general statement.

\subsection{Finite calibration}
\begin{table}[p]\footnotesize
\caption{Finite checks (exact integer or rational arithmetic unless stated).}\label{tab:finite}
\begin{tabular}{p{0.25\textwidth}p{0.35\textwidth}p{0.3\textwidth}}
\toprule
statement & instances & result\\
\midrule
Theorem~\ref{thm:eps} & $m=5$: natural, $h=0$, adversarial, three random $h$; $m=7$: two $p$ with natural, $h=0$, random $h$; $m=5$ with $\wt(\rho)=3,3,5$ & $\Delta(B)=q$ in every case; $\cZ$-bound $2q/2^{\wt(s)}$ never violated\\
Theorem~\ref{thm:W2} & $m=5$, author DDT, exact matrix products & $N\cdot\EDP_2$: natural $185/128$, $h=0$ $311/128$, adversarial $31/16$, random $h$ $303/256$; $x^{-1}$ with $d_0$ (ring-polynomial witness): $159/128$; $\wt(\rho)=3$ example $75/64$ ($\le6+6/q$)\\
Theorem~\ref{thm:UE} & $m=5$: 11 tables (incl.\ $\wt(\rho)=2,3$); $m=7$: 8 tables (incl.\ $\wt(\rho)=4$) & identity of Lemma~\ref{lem:Eid} holds for every $\cG$ direction; $\mathcal E_{\cZ}=211/512$ ($m=5$), $2059/8192$ ($m=7$); $\mathcal E_B\in[3.497,3.559]$ ($m=5$), $[3.345,3.357]$ ($m=7$)\\
Section~\ref{sec:repair} ($d_0$) & power sums $m=3,5,7,9$; coefficients: 40 cases each for $m=3,5$, 8 for $m=7$; correction rule: all $\lambda$, 800 pairs; $m=5$ all $1024$ local keys for $h=a^3$ & all agree; minimum number of non-zero $c$: $27$ ($\tau=0$, bound $18$), $26$ ($\tau=1$, bound $19$)\\
Lemma~\ref{lem:norm}, Theorems~\ref{thm:top}--\ref{thm:repair}, Proposition~\ref{prop:affine} (general $d$) & own arithmetic, $m=3,4,5$; 19 maps $d$ (random $\rho$, hyperplane and $K$, plus $d_0$); coefficients for 144/144/168 keys, all $L$ and $\mu$ for $m=3$ (45\,056 keys); $\chi$ two choices; 550 affine $h$ & all agree; minimal non-zero counts $6/13/28$ ($\tau\ne\eta A$) and $7/12/28$ ($\tau=\eta A$); Boolean degree of $g_d$ exactly $2$; pass rate in $\mathcal Q_2$ ($m=3$, sample) $0.873$--$0.879$\\
Proposition~\ref{prop:pf} & author DDT $\to$ integer Walsh, $m=3$ (24\,576) and $m=5$ (7\,864\,320 components), 16 tables, $\wt(\rho)=1,2,3$ & all agree; $\cG$ entries in $\{0,\pm4\}/q^2$ only\\
\bottomrule
\end{tabular}
\end{table}
\begin{table}[p]\footnotesize
\caption{Finite checks (continued).}\label{tab:finite2}
\begin{tabular}{p{0.25\textwidth}p{0.35\textwidth}p{0.3\textwidth}}
\toprule
statement & instances & result\\
\midrule
Theorem~\ref{thm:R1} & $m=3,5,7$ (20 tables; numerical norms) & every block below its bound; $\|M_\cZ\|^2=(q/2-1)/q^2$ exactly\\
Proposition~\ref{prop:Zexact} & independent code, $m=3,5$; $d_0$ and $1\oplus K_0(x/\rho)$ ($\rho=7,11$); $h\in\{0,a^3,\text{random}\}$; 15 tables & integer Walsh of the $\cZ$ rows $=x_s$; Gram $q^2[s=t]$ exactly; SVD norm and exact Rayleigh quotient of the witness $=(q/2-1)/q^2$\\
Theorem~\ref{thm:G}, Proposition~\ref{prop:Minv} & two independent implementations; $m=3$: 16, $m=5$: 21 tables; adversarial $h$ (hill climbing); $\mathsf M$ for $p=x^{-1}$, $m=3,\dots,11$ & $g=\Psi\Sigma$ entrywise; closed form of $\mathcal B^0_z$ exact; largest observed $\|M_\cG\|^2q/\mathsf M$ (floating point): $0.46287$ ($m=3$), $0.06390$ ($m=5$); $\mathsf M\le5$ for $d_0$ type\\
Theorem~\ref{thm:steps}, Corollary~\ref{cor:inv13} & every elementary inequality for $m=5,\dots,61$ (outward rational square roots); start lemma for $m=3$, all $4095$ two-word inputs, $M=[[1,1],[1,t]]$ & all hold; $\max\|\pi_2\|_\infty=0.00464<1/64$ ($q=8$ calibrates the start lemma only)\\
Example~\ref{ex:quant}, Proposition~\ref{prop:mc} & $m=3$: all $64$ carry-mask keys and all $512\cdot8$ linear $(L,\mu)$; $m=5$: $g_d$ for $\rho=1$ and $\rho=7$ & $\ell(h)=0$, no carry-mask key vanishes, $(L,\mu)=((3,0,0),0)$ vanishes ($64$ linear pairs in all); quadratic rank modulo affine $=4=m-1$, the $(m-1)$-AND circuit agrees at all points\\
Theorem~\ref{thm:V} & decomposition: toy $m=2$, $w=2$, all $16\,777\,216$ candidates, 32 cases; $m=3$ sample of $58\,035$ & accepted sets identical\\
Theorem~\ref{thm:V}(3) & $M=5$, $Q=3$, all $120$ permutations & $\RP$ law $=\RF\mid\mathsf D$\\
\bottomrule
\end{tabular}
\end{table}

\subsection{Formalization}
The Lean developments use Lean 4.31.0-rc1 and Mathlib revision
\texttt{d568c8c0}\allowbreak\texttt{9630de09}\allowbreak\texttt{7a046763}\allowbreak\texttt{c17b9ea9}\allowbreak\texttt{9f95f950}; every listed declaration depends only on the
standard axioms \texttt{propext}, \texttt{Classical.choice}, \texttt{Quot.sound}, with no \texttt{sorry}.
\begin{table}[!htbp]\footnotesize
\caption{What is checked in Lean, and what is an input.}\label{tab:lean}
\begin{tabular}{p{0.2\textwidth}p{0.40\textwidth}p{0.3\textwidth}}
\toprule
file & checked content & not covered / input\\
\midrule
\texttt{WarmStart} & arithmetic of the warm-started recursion: from initial values $\le x$ the bound $4x^2$; $\le4/q^2$ for real $q\ge32$ & the counts of Section~\ref{sec:diff} on a concrete ring\\
\texttt{SuffixReset}, \texttt{ResetCompose} & DDT row/column sums on any finite abelian group; reset identity of Lemma~\ref{lem:reset}; two applications (four layers) & the bound on the first part is an input\\
\texttt{W2Bridge} & averaged kernel of a parallel layer with product keys is the tensor product; real averaged DP of the $w=2$ superbox equals the abstract MDS kernel; $(2,2,1,1)$, $w=2$: $\E\DP\le b_\star^2$ & the one-word bound $b_\star$ (Theorem~\ref{thm:W2}) is an input; general $w$; identification with the reference code\\
\texttt{CarryPD} & two-digit addition with an \emph{arbitrary} carry function, $B=J\circ F$, real DDT: in $\cG$ directions ($d(a+u)=d(a)+k$, $k\ne0$) $\delta_B\le q$ for all $p,h,\beta$ & $\cZ$ and $\cE$ directions; Theorem~\ref{thm:W2}; connection to \texttt{W2Bridge}\\
\bottomrule
\end{tabular}
\end{table}
\begin{table}[!htbp]\footnotesize
\caption{What is checked in Lean (continued).}\label{tab:lean2}
\begin{tabular}{p{0.2\textwidth}p{0.40\textwidth}p{0.3\textwidth}}
\toprule
file & checked content & not covered / input\\
\midrule
\texttt{CarryZ}, \texttt{CarryE\_v2} & concrete carry $\star$ (product in a fixed basis $e:\Fq\simeq\FF^m$); fibres of $x\mapsto x\star s$ have $\le q/2$ points for $s\ne0$; $\cZ$ direction $\delta_B\le q$ ($d\ne0$, $d$ two-to-one); $\cE$ direction $\delta_B\le q$ ($p$ APN along $\rho$) & exact fibre size $q/2^{\wt s}$\\
\texttt{OneRound} & all $\alpha\ne0$: $\delta_B\le q$ and $\DP\le1/q$, for $d=\varepsilon\oplus K$, $\ker K=\{0,\rho\}$, $d\ne0$, $p$ APN, every $h$ (two-to-one derived; the $\cG$ lemma is that of \texttt{CarryPD}) & ---\\
\texttt{Transport\_v2} & for any group $G$ and bijection $\varphi:G\to\Fq^2$ carrying $+$ to the carry addition: $\mathrm{ddt}=\mathrm{ddtC}$ exactly; keyed families of PD layers have $\mathrm{avgDP}\le1/q$ in the definitions of \texttt{W2Bridge} & the identification $\GR(4,m)\cong(\Fq^2,\oplus_\star)$ (Lemma~\ref{lem:add}) is the hypothesis $\varphi$; the two-block bound $4/N$\\
\texttt{RepairCore\_v2} & $A^2+B\ne0\Rightarrow$ top coefficients not both $0$; root count $\ge q-(s-1)$; for any finite $\Fq$-module $V$ and $\ell(g)=1$: surjectivity, $\ell(h+\lambda g)=\ell(h)+\lambda$, exactly one failing $\lambda$, $q\cdot\#\ker\ell=|V|$ & $\ell=\ell_d$, power sums, Lucas, $\ell_d(g_d)=1$, $\Psi=\Phi^2$\\
\texttt{ColsB\_v3}, \texttt{ColsF\_v2} & the rows of Proposition~\ref{prop:table} as counts of the real DDT of $B$ and of $F$ on concrete carry coordinates, for $\wt(\rho)=1$: $\cZ\to\cZ$ $0$, $\cZ\to\cE$ column $\le q^2/2$, $\cE\to\cZ$, $\cE\to\cE$ columns $\le2q$, $\cG\to\cZ$ $\le2$, $\cG\to\cE$ $\le4$; final layer: low digit $0$ only from $\cZ$, $\cZ$ column $\le q^2$, $\cE$ column $\le2q$, $\cG$ entries $\le2$ & general $\omega$ (only $\le1/2$ for $\cZ\to\cE$), so Theorem~\ref{thm:W2}(2)\\
\texttt{TwoStep}, \texttt{Bridge} & Theorem~\ref{thm:S} one step from the warm start and the final layer: $(T_B^2T_F)(\alpha,\gamma)\le4/N$ ($q\ge12$); the real key average of $S_2$ (with whitening) equals $T_B^2T_F$; hence \texttt{hb1}--\texttt{hb4} of \texttt{W2Bridge} are theorems and the $(2,2,1,1)$, $w=2$ average is $\le(4/N)^2$ & $\wt(\rho)=1$ only; $t\ge2$ of Theorem~\ref{thm:S}; the identification $\varphi$ of $\GR(4,m)$ with $(\Fq^2,\oplus_\star)$ remains a hypothesis\\
\texttt{RepairLink\_v2} & $\ell_d$ on actual definitions ($A_d$, $B_d$ over a transversal, Frobenius) is $\Fq$-linear; $\Delta_{g_d}=r$; repair, one failure and pass count $1-1/q$ for the actual $\ell_d$ & $\ell_d(g_d)=1$ (power sums $p_2,p_3$) is the only hypothesis; that the index set is a transversal is not stated in Lean\\
\bottomrule
\end{tabular}
\end{table}
\begin{table}[!htbp]\footnotesize
\caption{What is checked in Lean (continued): the concrete two-word mixer.}\label{tab:lean3}
\begin{tabular}{p{0.2\textwidth}p{0.40\textwidth}p{0.3\textwidth}}
\toprule
file & checked content & not covered / input\\
\midrule
\texttt{DigitSplit\_v2} & the digit split $\varphi_4:(\mathbb Z/4)^m\to\Fq^2$, $x=a+2b$ coordinatewise, carries $+$ to the carry addition (\texttt{phi4\_add}) & the ring multiplication of $\GR(4,m)$ is not used; a statement there that takes the mixer entries in the \emph{pointwise} ring $(\mathbb Z/4)^m$ has no instance (\texttt{VacuousZ4})\\
\texttt{VacuousZ4} & in the pointwise ring $(\mathbb Z/4)^m$, $m\ge1$, four unit entries and a unit determinant are incompatible (\texttt{w2\_Z4\_mixer\_hyps\_false}); hence the earlier specialization of the two-word bound to that ring is vacuous and is not counted as a verification & ---\\
\texttt{MixAdd} & two-word MDS kernel lemma over any finite abelian group $A$ for an additive automorphism $M$ of $A^2$ whose four partial maps are bijective (\texttt{mds2A\_le}); real key average of the $w=2$ superbox $=$ this kernel; $(2,2,1,1)$, $w=2$: $\E\DP\le\varepsilon^2$ (\texttt{w2A\_split\_tail\_le}) & the one-word bound $\varepsilon$ is an input\\
\texttt{ThetaZ4\_v3} & mixer $M(x,y)=(x+y,\,x+\Theta y)$ on $A=(\mathbb Z/4)^m$ for an additive automorphism $\Theta$ with $\Theta-I$ invertible; with the actual local blocks $S_2$ and $\varphi_4$: $(2,2,1,1)$, $w=2$ average $\le(4/N)^2$ (\texttt{w2\_Z4\_theta\_le}); explicit witness $m=5$, $\Theta=$ multiplication by $t$ in $(\mathbb Z/4)[t]/(t^5+t^2+1)$, with $\Theta^{-1}$ and $(\Theta-I)^{-1}$ (\texttt{theta5}, \texttt{theta5\_mixer\_mds}, \texttt{w2\_Z4\_theta5\_le}) & PD data $(p,K,\varepsilon)$ with $\wt(\rho)=1$ are hypotheses; for general $m$ the choice of $\Theta$ is on paper (Remark below)\\
\bottomrule
\end{tabular}
\end{table}
The one-round bound $\varepsilon(B)\le1/q$ of Theorem~\ref{thm:eps} is formalized for all directions on concrete
carry coordinates. For $\wt(\rho)=1$ and $w=2$, the chain Proposition~\ref{prop:table} $\to$ Theorem~\ref{thm:S} (one step) $\to$
Theorem~\ref{thm:W2}(1) $\to$ Theorem~\ref{thm:P4}(1) is checked in Lean on the real key averages on the additive group
$(\mathbb Z/4)^m$ with the digit split $\varphi_4$, for the two-word mixer $(x,y)\mapsto(x+y,\,x+\Theta y)$
(\texttt{w2\_Z4\_theta\_le}), and the mixer hypotheses are shown to be satisfiable ($m=5$). The mixer is formalized as an additive
automorphism, not as a matrix of ring scalars: in the pointwise ring $(\mathbb Z/4)^m$ a matrix with unit entries and unit
determinant does not exist (\texttt{VacuousZ4}), so the earlier specialization with such scalars proves nothing about a concrete mixer
and is superseded. That the additive group of $\GR(4,m)$ in a polynomial basis is $(\mathbb Z/4)^m$, and that multiplication by $t$ is
the companion matrix $\Theta$ of the lift $\hat f$, are used on paper; the local bound for $\wt(\rho)>1$, general $\omega$, general $w$,
Theorem~\ref{thm:UE} and Section~\ref{sec:mixing} are not formalized; of Section~\ref{sec:repair} the linear part on the actual
definitions is formalized, the power sums are not. All Lean sources needed by these declarations are released as one source
closure that builds from sources only.

\begin{remark}[The mixer for general $m$]\label{rem:theta}
Let $f\in\FF[t]$ be irreducible of degree $m\ge2$ and $\hat f$ its lift with coefficients in $\{0,1\}$. Let $\Theta$ be the matrix of
multiplication by $t$ on $(\mathbb Z/4)[t]/(\hat f)\cong(\mathbb Z/4)^m$. Then $\det\Theta=\pm\hat f(0)$ and $\det(\Theta-I)=\pm\hat f(1)$
are odd, hence units of $\mathbb Z/4$, so $\Theta$ and $\Theta-I$ are invertible and $(x,y)\mapsto(x+y,\,x+\Theta y)$ is the mixer
$\bigl[\begin{smallmatrix}1&1\\1&t\end{smallmatrix}\bigr]$ of $\GR(4,m)$ in additive form, with partial maps $I,I,I,\Theta$.
This was checked exactly for all $69$ binary irreducible $f$ with $2\le\deg f\le8$ (inverse identities), and the agreement with
multiplication by $t$ in $\GR(4,m)$ on all pairs for $m=3,5$ (two polynomials each).
\end{remark}

\subsection{Cost}
Counts are in a bit-level model with AND gates (XOR and NOT counted separately), for $m=5$. They are
reported to locate the constructions, not to claim an advantage.
\begin{table}[ht]\small
\caption{Bit-model counts ($m=5$).}\label{tab:cost}
\begin{tabular}{p{0.36\textwidth}p{0.14\textwidth}p{0.4\textwidth}}
\toprule
object & AND & averaged differential bound\\
\midrule
three-round Feistel with $x^3$, independent XOR keys & 75 & $4/N$ (\cite{AO}, Thm.~2; bijective round function, $\ge3$ rounds)\\
$S_2$ with the natural table of \cite{FukuiFFA} & 145 & $4/N$ (Theorem~\ref{thm:W2})\\
$S_2$ with an inverse-based PD table & 235 & $4/N$ (Theorem~\ref{thm:W2})\\
one $F$, PD table with $h=a^3$ & 75 (98 XOR, 1 NOT) & ---\\
correction $g_d$ alone (Proposition~\ref{prop:mc}) & $m-1$, minimal (4 for $m=5$) & --- (cost of the correction only, not of $h$)\\
\bottomrule
\end{tabular}
\end{table}
The Feistel network reaches the same constant with fewer AND gates. The differences that remain are
structural (ring differences with carry versus XOR differences; $2m$-bit ring key additions; direction-wise
structure; uniformity over all $h$; the integral structure of Section~\ref{sec:repair}), not an ordering by
cost. Ring-only and MPC costs were not measured.
\clearpage

% sec_8_limits_v4.tex -- C9 paper lane (P9-01): exponent 1/2 stated for d=1+K0(x/rho) with the general-d constant;
% otherwise unchanged from sec_8_limits_v3.tex.
% sec_8_limits_v3.tex -- C8: exponents, step count, quantifiers of Section 4, formalization status.
% sec_8_limits_v2.tex -- main author (C7): general-d scope, R1 status.
\section{Limits and remaining design questions}\label{sec:limits}

\subsection{What the theorems cover}
Each theorem fixes its quantifiers as follows.
Theorems~\ref{thm:eps} and~\ref{thm:UE}: all $h$, all $\rho$, all $(p,d)$ under Assumption~\ref{ass:PD} (for
Theorem~\ref{thm:UE} the permutation property of $p$ is not used), exact counts.
Theorem~\ref{thm:W2}: odd $m\ge5$, all $h$; constant $4$ only for $\wt(\rho)=1$ in the fixed basis,
$(6+6/q)$ otherwise; average over independent uniform keys and whitening.
Theorem~\ref{thm:P4}: the $(2,2,1,1)$ arrangement with three mixings and reset keys $L_2,L_3$,
$M_1$ MDS, $M_2,M_3$ invertible, average over $10w$ independent uniform key words.
Theorem~\ref{thm:V}: non-adaptive queries on pre-fixed disjoint slices; the real construction is accepted for
every key; false acceptance bounded by a union over $N^{3w}$ candidates.
Section~\ref{sec:repair}: $p(x)=x^{q-2}$, every $d$ satisfying (P2), every $m\ge3$ for the algebra (odd $m\ge5$ when
combined with Section~\ref{sec:diff}), every local key; the certificate is sufficient for every $\FF$-linear $L$, and its
necessity is not asserted (Example~\ref{ex:quant}); the degree statement concerns $h$ only; Proposition~\ref{prop:mc} concerns $g_d$ alone.
Theorem~\ref{thm:R1}, Proposition~\ref{prop:Zexact}, Theorem~\ref{thm:G}: every $m\ge3$, all PD data, all $h$; averaged transitions.
Corollary~\ref{cor:half}, Corollary~\ref{cor:inv13}: $p(x)=x^{q-2}$, odd $m$, $d=1\oplus K_0(x/\rho)$, all $h$.
Theorem~\ref{thm:steps}: $m\ge5$, all PD data, all $h$, the alternating arrangement A1 with independent uniform keys in every step and a
two-word mixing with unit entries and determinant.

\subsection{What remains open or is not addressed}
\begin{itemize}
\item The one-word contraction has exponent $\gamma\approx0.2075$ for all PD data (Theorem~\ref{thm:R1}) and the optimal exponent
$1/2$ for the inverse map with $m$ odd and $d=1\oplus K_0(x/\rho)$ (Corollary~\ref{cor:half}). For the inverse map with another admissible
$d=\varepsilon\oplus K$, $K$ of linearized degree $e$, the constant becomes $2^{e+1}+5$ (Proposition~\ref{prop:Minv}), which is not uniform in $m$ when $e$
grows with $m$. Exponent $1/2$ uniformly over all PD data is reduced to the
$h$-free bound $\max_z\|\mathcal B^0_z\|\le Cq^2$ and is open. The step counts (Theorem~\ref{thm:steps}, Corollary~\ref{cor:inv13}) have
large constants; they are statements about the scale $1/N$, not round recommendations. Numerical contraction constants are estimates.
\item No fixed-key differential bound, no linear or algebraic bounds, no related-key statement.
\item No key schedule; the $10w$ independent key words are the randomness of the averaging model.
\item No lower bound on attack time; the implemented enumerations are upper bounds for one method.
Meet-in-the-middle and character-sum (Fourier) decompositions of the candidate search, higher-order
integrals and division-property analyses were not carried out.
\item The local correction removes one family of local zero sums inside an $S_2$ block; it does not remove
the word-level zero sum used by Theorem~\ref{thm:V}, and it is not a statement about integral resistance of
the arrangement.
\item The classification of $h$ with $\ell_d(h)=0$ is not attempted; the vanishing example of
Example~\ref{dz:prop:X21} is a single instance. Non-inverse $p$ and non-affine $d$ are not treated in
Section~\ref{sec:repair}.
\item Cost: no claim of minimal AND count or of an advantage over known constructions; ring-only and MPC
costs are unmeasured.
\end{itemize}

\subsection{Design questions left for later work}
The arrangement $(2,2,1,1)$ and the number of mixing layers are fixed in this paper; we do not add layers
to raise attack cost. A design that addresses the word-level integral of Lemma~\ref{lem:V2}, a key
schedule, and fixed-key statements would be required before any construction built from these
components could be considered as a cipher.


\section*{Data and code availability}
The finite checks of Section~\ref{sec:verif} and the Lean~4 sources with their build records are provided in the public
repository \url{https://github.com/hirokifukui/residue-apn-carry}, archived at Zenodo,
\url{https://doi.org/10.5281/zenodo.23330906} (all versions).
The directory \texttt{checks/} contains three independently written checkers, with their expected outputs, for
Proposition~\ref{prop:Zexact}, Theorem~\ref{thm:steps} with Corollary~\ref{cor:inv13}, and Example~\ref{ex:quant} with
Proposition~\ref{prop:mc} (including the gate count of the correction $g_d$ in Table~\ref{tab:cost}). For every other row of
Tables~\ref{tab:finite}, \ref{tab:finite2} and~\ref{tab:cost}, and for the check reported in Remark~\ref{rem:theta}, the
directory \texttt{evidence/} contains the unedited scripts that produced the reported numbers, the logs, exit codes and result
files of those runs, and the input tables (included, or regenerated by included scripts); the index
\texttt{EVIDENCE\_COVERAGE.tsv} lists, row by row, the script, the recorded output and whether the row is re-run. The script
\texttt{run\_all\_representatives.sh} re-runs one representative command per row (Python~3 with NumPy and SymPy; about ten
minutes; for the gate count of one $F$ with $h=a^3$ through a small added wrapper that leaves out a cross-check against a
table that is not distributed; for the adversarial search of Theorem~\ref{thm:G} with $m=3$ by recomputing the $40$ recorded
witnesses, re-running the search itself being optional) and compares the result with the recorded output. For the following entries the numbers rest on the recorded outputs
only, because the runs are long: the $m=7$ instances of Theorems~\ref{thm:eps} and~\ref{thm:R1}, the $m=5$ instances of
Theorem~\ref{thm:G} and Proposition~\ref{prop:Minv}, and 29 of the 32 toy cases and the $m=3$ sample of Theorem~\ref{thm:V}.
A finite check is a calibration, not a proof.

\section*{Use of generative AI}
Generative AI systems, including Anthropic Claude and OpenAI GPT models, were used throughout this project to assist with
the exploration and comparison of proof strategies and derivations, adversarial review of successive drafts, the design and
implementation of the checker scripts for the finite computations, tactic-level search and code drafting for the Lean
formalization, and drafting and revision of the manuscript. The author made all final mathematical and editorial decisions,
checked every definition, statement, proof, and citation included in the manuscript, and assumes sole responsibility for the
entire content. As described in Section~\ref{sec:verif}, designated finite computations and selected proof components were
additionally checked by checker scripts or by the Lean kernel; these checks do not rely on accepting AI-generated
mathematical assertions, and a finite check is a calibration, not a proof.

\appendix
% app_A_mixing_v2.tex = app_A_mixing_v1 with one inline sum of the proof of Lemma app-line set as a display
% (anonymous iacrtrans build: 24.5pt overfull box); no change of content.
% app_A_mixing_v1.tex -- C8 (P8-04, GPT C7-F02): the C5/C6 lemmas used by Section 6, collected and proved here so that
% Section 6 depends on no research note. Sources (unchanged originals): c5/a1/notes/LYAP_P5.md (L1-L5, Prop. 1),
% c6/a1/notes/LYAP_P6.md (corrected meaning of Lambda', L4', Prop. 1', comparison matrix), c6/a1/notes/N1_P6.md
% (W1-W4, G0, G1, Proposition D2). Not included because Section 6 does not use them: Lambda'' and the axis low modes (LYAP_P6 s.2),
% the zeta_2 reduction and R2 (N1_P6 s.7), Perron weights (LYAP_P6 s.5).
\section{The two-word reduction}\label{app:mix}

Throughout, Assumption~\ref{ass:PD} holds, $T=T_B$, $W=R\setminus\{0\}$, and $M$ is a $2\times2$ matrix over $R$ whose entries and
determinant are units. All spaces carry the counting-measure inner product.

\subsection{Spaces and operators}
Let $V_0=\{X:R^2\to\mathbb R,\ X(0,0)=0\}$, $U=\mathbf 1_{R^2\setminus\{(0,0)\}}/(N^2-1)$, and work in $V_0\ominus U$. For $X\in V_0$ put
$SX=T^\top XT$ (that is, $(SX)(\gamma_1,\gamma_2)=\sum_{\beta}T(\beta_1,\gamma_1)X(\beta_1,\beta_2)T(\beta_2,\gamma_2)$) and $\mathcal TX=\mathrm{Perm}_M(SX)$ with
$(\mathrm{Perm}_MX)(M\beta)=X(\beta)$. Since $T^\top$ is doubly stochastic with $T(0,\cdot)=e_0$ and $T(\alpha,0)=0$ for $\alpha\ne0$, and $\mathrm{Perm}_M$ is an
orthogonal permutation fixing $(0,0)$, $\mathcal T$ maps $V_0$ to itself, $\mathcal TU=U$, $\|\mathcal T\|\le1$, and $\mathcal T$ preserves $V_0\ominus U$.

\emph{Cells.} For $s,t\in\Fq$ let $f_s=\{\beta\in R:\mathrm{lo}\,\beta=s\}$ ($q$ points). The cells of $R^2\setminus\{(0,0)\}$ are $f_s\times f_t$ for
$(s,t)\ne(0,0)$ and $(f_0\times f_0)\setminus\{(0,0)\}$; every cell has at least $q^2-1$ points. $\Lambda'\subset V_0$ is the space of functions
constant on cells, $P=\Pi_{\Lambda'}$ and $Q=I-P$ (on $V_0\ominus U$; note $U\in\Lambda'$). In one word, $\lambda$ is the space of functions on
$R$ constant on the fibres $f_s$, $\tilde\lambda$ the functions on $W$ constant on $f_s\cap W$, and $\perp\subset\ell^2(W)$ the functions with sum $0$
on every $f_s\cap W$; $T'=T|_{W\times W}$ and $\tau=\|\Pi_\perp T'^\top\Pi_\perp\|$ as in Theorem~\ref{thm:R1}.

\subsection{The slow mode}
\begin{lemma}[L1]\label{lem:L1}
$\mathrm{Perm}_M$ maps cells to cells and fixes $(0,0)$; hence $\mathrm{Perm}_MP=P\,\mathrm{Perm}_M$.
\end{lemma}
\begin{proof}
Reduction modulo $2R$ is a ring homomorphism, so $\mathrm{lo}(M\beta)=\bar M\,\mathrm{lo}(\beta)$ with $\bar M$ invertible over $\Fq$.
\end{proof}

\begin{lemma}[L2]\label{lem:L2}
In one word, $\Pi_\lambda T^\top\Pi_\lambda$ is the projection onto the constants: the low digit of the output, averaged over an input fibre
$f_u$, is uniform.
\end{lemma}
\begin{proof}
Use Lemma~\ref{lem:diffB}. $\cG$ ($u\notin\{0,\rho\}$): for each $(a,s)$ exactly one $b$ gives low difference $s$, so every row of the fibre has a uniform
low marginal. $\cZ$ ($u=0$): $v=0$ is the zero difference (output $0$), and for $v\ne0$ the low output $s=d(a)v$ is uniform on $v\,\Img d$ (each value
with probability $2/q$); averaging over $v$ gives $q^{-1}([s=0]+[s\ne0]\cdot(2/q)\cdot\#\{v\ne0:s\in v\Img d\})=q^{-1}([s=0]+[s\ne0])=1/q$.
$\cE$ ($u=\rho$): the low output $s_v(a)$ is a bijective function of $v$ for each $a$, so the average over $v$ is uniform.
\end{proof}

\begin{lemma}[L3]\label{lem:L3}
$\|\Pi_\lambda T^\top\Pi_{\lambda^\perp}\|\le\sqrt{2/q}$.
\end{lemma}
\begin{proof}
For $Y\in\lambda^\perp$ (sum $0$ on every input fibre) the low marginal of $T^\top Y$ is $m(s)=\sum_\alpha Y(\alpha)K(\alpha,s)$ with
$K(\alpha,s)=\Pr[\mathrm{lo}\text{ of output}=s\mid\alpha]$, and $\|\Pi_\lambda T^\top Y\|^2=q^{-1}\sum_sm(s)^2$. Rows in $\cG$ have constant $K$ (Lemma~\ref{lem:L2}),
so their contribution vanishes. On the $\cZ$ fibre, $K_ZK_Z^\top=\mathrm{diag}(1)\oplus q^{-1}(I+J)$, because $|v\Img d\cap v'\Img d|=q/4$ for $v\ne v'$ ($\Img d$ is
an affine hyperplane); on vectors of sum $0$ its largest Rayleigh quotient is $1$, so $\|K_Z^\top Y_Z\|\le\|Y_Z\|$. On the $\cE$ fibre
the $q\times q$ matrix $K_E$ is doubly stochastic, so $\|K_E^\top Y_E\|\le\|Y_E\|$. Hence $\|m\|\le\|Y_Z\|+\|Y_E\|\le\sqrt2\|Y\|$ and
$\|\Pi_\lambda T^\top Y\|=q^{-1/2}\|m\|\le\sqrt{2/q}\,\|Y\|$.
\end{proof}

\begin{lemma}[L4]\label{lem:L4}
On $V_0\ominus U$, $P\mathcal TP=-q^{-2}\psi\psi^\top/\|\psi\|^2$ with $\psi=-\mathbf 1_c+q^{-2}\mathbf 1_{R'}$, where $c=(f_0\times f_0)\setminus\{(0,0)\}$ and
$R'$ is the set of the other non-zero states; in particular $\|P\mathcal TP\|=1/N$.
\end{lemma}
\begin{proof}
For a cell $\sigma=(s,t)\ne(0,0)$, $S\mathbf 1_{f_s\times f_t}=(T^\top\mathbf 1_{f_s})\otimes(T^\top\mathbf 1_{f_t})$; by Lemma~\ref{lem:L2} each factor has mass $1$ in every
fibre and value $[s=0]$ at $0$. Hence the cell projection of $S\mathbf1_\sigma$ equals $q^{-2}$ on $R'$ and $1/(q^2-1)$ on $c$, while
$S(\mathbf 1_{f_0\times f_0}-e_{00})$ projects to $q^{-2}$ on $R'$ and $0$ on $c$. For $X=\sum_\sigma x_\sigma\mathbf 1_\sigma\in\Lambda'$ of mass $0$,
$q^2\sum_{\sigma\ne00}x_\sigma=-(q^2-1)x_{00}$, and substituting gives $P SX=(x_{00}/q^2)\psi$; $\mathrm{Perm}_M$ preserves $c$ and $R'$ (Lemma~\ref{lem:L1}). Finally
$\langle\psi,X\rangle/\|\psi\|^2=-x_{00}$.
\end{proof}

\begin{lemma}[L5]\label{lem:L5}
$\eta:=\|P\mathcal TQ\|\le\eta_{\mathrm{rss}}=\sqrt{2/q+1/(q^2-1)}$.
\end{lemma}
\begin{proof}
Work in $\ell^2(R^2)\supset V_0$. There, the cell-constant functions together with $e_{00}$ form $(\lambda\otimes\lambda)\oplus\mathrm{span}\{e'\}$,
$e'=e_{00}-q^{-2}\mathbf 1_{f_0\times f_0}$, an orthogonal sum with $\|e'\|^2=1-q^{-2}$; for $Y\in V_0$ the projection onto $\Lambda'$ is the projection onto this
space, because $Y\perp e_{00}$. By Lemma~\ref{lem:L1} $\mathrm{Perm}_M$ may be removed. For $X\in\Lambda'^\perp\subset(\lambda\otimes\lambda)^\perp$ split $X$ into its $\lambda\otimes\lambda^\perp$, $\lambda^\perp\otimes\lambda$ and
$\lambda^\perp\otimes\lambda^\perp$ parts; $\Pi_{\lambda\otimes\lambda}S=(\Pi_\lambda T^\top)\otimes(\Pi_\lambda T^\top)$ and Lemma~\ref{lem:L3} bound the three images, which are mutually
orthogonal, by $\sqrt{2/q}$ times the norms, so the $\lambda\otimes\lambda$ component is at most $\sqrt{2/q}\|X\|$. The $e'$ component is
$|\langle(T\otimes T)e',X\rangle|/\|e'\|=q^{-2}|\langle(T\mathbf1_{f_0})\otimes(T\mathbf1_{f_0}),X\rangle|/\|e'\|\le q^{-2}\cdot q\|X\|/\sqrt{1-q^{-2}}$ (the vector
$(T\mathbf 1_{f_0})\otimes(T\mathbf1_{f_0})$ has norm at most $q$). The two components are orthogonal.
\end{proof}

\subsection{One word, and the mixed Gram estimates}
\begin{lemma}[W1--W4]\label{lem:W}
(W2) $\ell_W:=\|\Pi_{\tilde\lambda}T'^\top\Pi_\perp\|\le\sqrt{2/(q-1)}$. (W3) $\nu:=\|T'^\top\Pi_\perp\|$ satisfies $\nu^2\le2/(q-1)+\tau^2$; hence, with
$\mathcal K=\perp\otimes\perp$ and $\mathcal R=(\tilde\lambda\otimes\perp)\oplus(\perp\otimes\tilde\lambda)$, $\|S\Pi_{\mathcal K}\|\le\nu^2$ and $\|S\Pi_{\mathcal R\oplus\mathcal K}\|\le\sqrt2\nu$.
(W4) $S$ preserves the axis $\mathrm{Ax}=e_0\otimes\ell^2(W)\oplus\ell^2(W)\otimes e_0$ and the interior $\ell^2(W\times W)$, and $\|S|_{\mathrm{Ax}}\|\le1$.
(W1, not used below) $\|\Pi_{\tilde\lambda\ominus1}T'^\top\Pi_{\tilde\lambda\ominus1}\|=1/q$.
\end{lemma}
\begin{proof}
(W2) Extend $Y\in\perp$ by $0$ at the origin; it lies in $\lambda^\perp$ on $R$ and $(T^\top Y)(0)=Y(0)=0$. Averaging over $q-1$ instead of $q$ points in $f_0$
multiplies the bound of Lemma~\ref{lem:L3} by at most $\sqrt{q/(q-1)}$. (W3) $T'^\top$ preserves mass, so $T'^\top Y$ has no constant component for $Y\in\perp$;
split it into its $\tilde\lambda$ and $\perp$ parts. On the interior $S=T'^\top\otimes T'^\top$, so $\|S\Pi_{\mathcal K}\|=\nu^2$; writing
$\mathcal R\oplus\mathcal K=(\ell^2(W)\otimes\perp)\oplus(\perp\otimes\tilde\lambda)$ (orthogonal inputs, each image of norm at most $\nu$) gives $\sqrt2\nu$. (W4) $T^\top e_0=e_0$ and
$T^\top\ell^2(W)\subset\ell^2(W)$. (W1) is a rank-one computation from Lemma~\ref{lem:L2} and is recorded only for completeness.
\end{proof}

\begin{lemma}[G0, G1]\label{lem:G}
(G0) If $X\perp\Lambda'$, then $\|\Pi_{\tilde\lambda\otimes\tilde\lambda}X_{\mathrm{int}}\|\le\sqrt{2/(q-1)}\,\|X_{\mathrm{Ax}}\|$.
(G1) $\|(I-\Pi_{\mathcal K})\,\mathrm{Perm}_M\,\Pi_{\mathrm{Ax}}\|\le2/\sqrt{q-1}$.
\end{lemma}
\begin{proof}
(G0) On cells $(s,t)$ with $s,t\ne0$ the sum of $X$ is $0$ and there are no axis points, so the $\tilde\lambda\otimes\tilde\lambda$ component vanishes. A cell
$(0,t)$, $t\ne0$, has $q$ axis points and $q(q-1)$ interior points; the interior sum is minus the axis sum, which is at most $\sqrt q$ times the axis norm,
so the component has norm at most $\|X_{\mathrm{Ax}}|_{\mathrm{cell}}\|/\sqrt{q-1}$. The cell $c$ contributes in the same way with $2(q-1)$ axis points.
(G1) $\mathrm{Perm}_M(e_0\otimes g)$ is supported on the line $\{(m_{12}r,m_{22}r)\}$ (both coefficients units), which lies in the interior; for fixed second
coordinate at most one first coordinate in a fibre lies on the line, and conversely. Hence each of the two one-word $\tilde\lambda$ projections of the
line function has norm at most $\|g\|/\sqrt{q-1}$; as $I-\Pi_{\mathcal K}=\Pi_{\tilde\lambda}\otimes I+\Pi_\perp\otimes\Pi_{\tilde\lambda}$ is an orthogonal sum, the
non-$\mathcal K$ part is at most $\sqrt{2/(q-1)}\|g\|$. The two halves of the axis are orthogonal inputs, which gives $\sqrt2\cdot\sqrt{2/(q-1)}=2/\sqrt{q-1}$.
\end{proof}

\begin{proposition}[D2]\label{prop:appD2}
\[
\theta_2=\|Q\mathcal T^2Q\|\le\eta_{\mathrm{rss}}+\sqrt{\bigl(\nu^2+(2+\sqrt2)/\sqrt{q-1}\bigr)^2+2\nu^2}.
\]
\end{proposition}
\begin{proof}
With $A=P\mathcal TP$, $B=P\mathcal TQ$, $C=Q\mathcal TP$, $D=Q\mathcal TQ$ one has $Q\mathcal T^2Q=CB+D^2$, $\|C\|\le1$ and $\|B\|\le\eta_{\mathrm{rss}}$ (Lemma~\ref{lem:L5}).
Let $X=QX$, $a=\|X_{\mathrm{Ax}}\|$, $b=\|X_{\mathrm{int}}\|$. Then $SX=SX_{\mathrm{Ax}}+S\Pi_{\tilde\lambda\otimes\tilde\lambda}X_{\mathrm{int}}+S\Pi_{\mathcal R\oplus\mathcal K}X_{\mathrm{int}}$, the last two
terms of norm at most $\sqrt{2/(q-1)}\,a$ (Lemma~\ref{lem:G}, G0) and $\sqrt2\nu b$ (Lemma~\ref{lem:W}). Write $Y=\mathrm{Perm}_MSX=L+E$ with
$L=\mathrm{Perm}_MSX_{\mathrm{Ax}}$, $\|L\|\le a$, and $\|E\|\le\sqrt{2/(q-1)}\,a+\sqrt2\nu b$. Since $\mathcal K\perp\Lambda'$, $QY=\Pi_{\mathcal K}L+Q(I-\Pi_{\mathcal K})L+QE$, and
$D^2X=Q\,\mathrm{Perm}_MS\,QY$ with $\|S\Pi_{\mathcal K}L\|\le\nu^2a$ and $\|(I-\Pi_{\mathcal K})L\|\le(2/\sqrt{q-1})a$ (G1). Collecting,
$\|D^2X\|\le(\nu^2+(2+\sqrt2)/\sqrt{q-1})a+\sqrt2\nu b$, and Cauchy--Schwarz with $a^2+b^2=\|X\|^2$ gives the bound for $\|D^2\|$.
\end{proof}

\subsection{Comparison, start and the line lemma}
\begin{lemma}[comparison]\label{lem:app-cmp}
For $a_j=\|PD_j\|$, $b_j=\|QD_j\|$: $a_{j+1}\le q^{-2}a_j+\eta b_j$, $b_{j+1}\le a_j+b_j$, hence
$a_{j+2}\le(\eta+q^{-4})a_j+\eta(1+q^{-2})b_j$ and $b_{j+2}\le\zeta_2a_j+\theta_2b_j$ with $\zeta_2=\|Q\mathcal T^2P\|\le1$.
\end{lemma}
\begin{proof}
$PD_{j+1}=P\mathcal TPD_j+P\mathcal TQD_j$ with Lemmas~\ref{lem:L4} and~\ref{lem:L5}; $\|QD_{j+1}\|\le\|D_j\|\le a_j+b_j$; iterate once for $a_{j+2}$, and split
$D_j=PD_j+QD_j$ under $Q\mathcal T^2$ for $b_{j+2}$.
\end{proof}

\begin{lemma}[start]\label{lem:app-start}
$\|D_1\|_2\le\|\pi_1\|_2\le q^{-1/2}$ and $\|\pi_2\|_\infty\le1/N$.
\end{lemma}
\begin{proof}
$\|\pi_1\|_2^2=\mathcal E(\alpha_1)\mathcal E(\alpha_2)$ with the row energies $\mathcal E(\alpha)=\sum_\beta T(\alpha,\beta)^2$, which equal $1$ for $\alpha=0$ and are at most
$\max_\beta T(\alpha,\beta)\le1/q$ otherwise (Theorem~\ref{thm:eps}); at least one word of $\alpha$ is non-zero. The second statement is Lemma~\ref{lem:D1}
($w=2$, $b_\star=1/q$) followed by the permutation $\mathrm{Perm}_M$.
\end{proof}

\begin{lemma}[line lemma]\label{lem:app-line}
For $m\ge5$ and every $j\in\{1,2\}$, $|\Pr_{\pi_{r+1}}[\beta_j=0]-1/(N+1)|\le(1779/256)\|D_r\|_2$.
\end{lemma}
\begin{proof}
Let $A_j=\{\beta\ne(0,0):\beta_j=0\}$ and $f=\mathcal T^*\mathbf 1_{A_j}$, so
\[f(\alpha)=\sum_\delta T(\alpha_1,\delta_1)T(\alpha_2,\delta_2)[(M\delta)_j=0].\]
The condition
$(M\delta)_j=0$ is the graph $\delta_2=\phi(\delta_1)$ of a bijection ($m_{j2}$ is a unit), so by Cauchy--Schwarz $f(\alpha)\le(\mathcal E(\alpha_1)\mathcal E(\alpha_2))^{1/2}$ and
$\|f\|_2\le\mathcal E_B\le1779/256$ (Theorem~\ref{thm:UE}). As $\mathcal TU=U$ and $D_r\perp U$,
$\Pr_{\pi_{r+1}}[\beta_j=0]=\langle f,U\rangle+\langle f,D_r\rangle$ with $\langle f,U\rangle=|A_j|/(N^2-1)=1/(N+1)$.
\end{proof}

\begin{lemma}[arithmetic]\label{lem:app-arith}
Let $m\ge5$, $x=(\sqrt3/2)^m$, $u=q^{-1/2}$. Then $x\le1/2$, $u/x\le3/8$, $u\le x^2$, $\eta_{\mathrm{rss}}\le\frac32u$, $2/(q-1)\le x^2/2$,
$(2+\sqrt2)/\sqrt{q-1}\le\frac{18}5u$; with Theorem~\ref{thm:R1}, $\nu^2\le\frac{11}2x^2$ and $\theta_2\le\frac{1919}{240}x<8x$. Moreover
$\eta_{\mathrm{rss}}+q^{-4}\le2u$ and $\eta_{\mathrm{rss}}(1+q^{-2})\le2u$.
\end{lemma}
\begin{proof}
$x\le1/2$ and $u/x=(2/3)^{m/2}\le3/8$ hold at $m=5$ and are monotone in $m$; $u\le x^2$ is $2^{3/2}\le3$. $\eta_{\mathrm{rss}}\le\frac32u$ is $q^2-4q-1\ge0$.
$2/(q-1)\le x^2/2$ is $4q^2\le(q-1)3^m$, true at $m=5$ ($4096\le7533$) and the ratio decreases. $(2+\sqrt2)^2q/(q-1)\le(7/2)^2\cdot32/31<324/25$.
Theorem~\ref{thm:R1} gives $\tau^2\le5x^2$, so $\nu^2\le\frac{11}2x^2\le\frac{11}4x$ and $\sqrt2\nu\le\sqrt{11}\,x\le\frac{10}3x$. Splitting the square root in
\eqref{eq:D2} by the triangle inequality, $\theta_2\le\nu^2+\sqrt2\nu+(\frac{18}5+\frac32)u\le(\frac{11}4+\frac{10}3+\frac{51}{10}\cdot\frac38)x=\frac{1919}{240}x$. The last two
inequalities follow from $\eta_{\mathrm{rss}}\le\frac32u$, $q^{-4}\le u/2$ and $\frac32(1+q^{-2})\le2$.
\end{proof}
% Exact rational re-check of every inequality for m=5..61: c8/checks/p803_steps.py part A (exit 0).


% bib_v2.tex = bib_v1 with [FukuiFFA] pointing to the Zenodo PAPER record of that preprint, version 1.0.1
% (10.5281/zenodo.23233176; concept 10.5281/zenodo.23232182) instead of the software record of that work. [FukuiIT] keeps 10.5281/zenodo.23053113 = concept DOI of the paper record of that preprint
% (all versions; the version archived at submission is 10.5281/zenodo.23074611).
% bib_v1.tex -- C9 paper lane (P9-14): shared bibliography of the public and the anonymous build.
% Changes from the bibliography of main_v6.tex: [BV] pages 65--81 and DOI (publisher page read 2026-10-12);
% uncited items [NK], [Nyberg93], [Square] removed; [EHRZ] DOI (OpenAlex record W7093092003); [FukuiIT], [FukuiFFA] cited as preprints with the Zenodo DOIs printed in the
% authors' own submitted manuscripts, without manuscript numbers (see c9/paper/AUTHOR_CHECK_C9.md, items B1-B3).
\begin{thebibliography}{99}

\bibitem{FukuiIT}
H.~Fukui, \emph{Optimal reduced representatives of quadratic APN permutations in odd dimension},
preprint, version 1.0.3, 2026, Zenodo, \url{https://doi.org/10.5281/zenodo.23053113}; submitted to IEEE Trans.\ Inform.\ Theory.

\bibitem{FukuiFFA}
H.~Fukui, \emph{Explicit reduced APN permutations: ramification and optimal Galois-ring lifts},
preprint, version 1.0.1, 2026, Zenodo, \url{https://doi.org/10.5281/zenodo.23233176}; submitted to Finite Fields Appl.

\bibitem{LMM}
X.~Lai, J.~L.~Massey, S.~Murphy, \emph{Markov ciphers and differential cryptanalysis},
Advances in Cryptology -- EUROCRYPT '91, LNCS 547, Springer, 1991, 17--38.

\bibitem{AO}
K.~Aoki, K.~Ohta, \emph{Strict evaluation of the maximum average of differential probability and
the maximum average of linear probability}, IEICE Trans.\ Fundamentals \textbf{E80-A} (1997), no.~1, 2--8.

\bibitem{Hong}
S.~Hong, S.~Lee, J.~Lim, J.~Sung, D.~Cheon, I.~Cho, \emph{Provable security against differential
and linear cryptanalysis for the SPN structure}, FSE 2000, LNCS 1978, Springer, 2001, 273--283.

\bibitem{Park}
S.~Park, S.~H.~Sung, S.~Chee, E-J.~Yoon, J.~Lim, \emph{On the security of Rijndael-like structures
against differential and linear cryptanalysis}, ASIACRYPT 2002, LNCS 2501, Springer, 2002, 176--191.

\bibitem{KW}
L.~R.~Knudsen, D.~Wagner, \emph{Integral cryptanalysis}, FSE 2002, LNCS 2365, Springer, 2002, 112--127.

\bibitem{Bernstein}
D.~J.~Bernstein, \emph{Stronger security bounds for permutations}, preprint, 2005-03-23,
\texttt{cr.yp.to/antiforgery/permutations-20050323.pdf}.

\bibitem{EHRZ}
A.~Ebeling, X.-D.~Hou, A.~Rydell, S.~Zhao, \emph{On sum-free functions},
Finite Fields Appl.\ \textbf{110} (2026), 102744, \url{https://doi.org/10.1016/j.ffa.2025.102744} (arXiv:2410.10426).

\bibitem{BV}
T.~Baign\`eres, S.~Vaudenay, \emph{Proving the security of AES substitution-permutation network}, Selected Areas in
Cryptography -- SAC 2005, LNCS 3897, Springer, 2006, 65--81, \url{https://doi.org/10.1007/11693383_5}.

\bibitem{HZ}
X.-D.~Hou, S.~Zhao, \emph{Further results on sum-freedom of binary and $q$-ary functions}, arXiv:2609.31489v1, 2026.

\end{thebibliography}


\end{document}
